% Vmem-phi -- Full draft sections: Corruption suite, Corruption algorithms, Theory, Hardware
% Companion outline: Docs/paper_skeleton.md ; measured numbers: Docs/paper_figures.md
%
% DUAL USE:
%   * Standalone: compiles as-is on Overleaf (the wrapper below supplies the
%     preamble, \begin{document}, and a manual bibliography).
%   * \input into a parent manuscript: define \theoryhardwareEMBED before the
%     \input and the standalone wrapper + bibliography are skipped, e.g.
%       \newcommand{\theoryhardwareEMBED}{}
%       % Vmem-phi -- Full draft sections: Corruption suite, Corruption algorithms, Theory, Hardware
% Companion outline: Docs/paper_skeleton.md ; measured numbers: Docs/paper_figures.md
%
% DUAL USE:
%   * Standalone: compiles as-is on Overleaf (the wrapper below supplies the
%     preamble, \begin{document}, and a manual bibliography).
%   * \input into a parent manuscript: define \theoryhardwareEMBED before the
%     \input and the standalone wrapper + bibliography are skipped, e.g.
%       \newcommand{\theoryhardwareEMBED}{}
%       % Vmem-phi -- Full draft sections: Corruption suite, Corruption algorithms, Theory, Hardware
% Companion outline: Docs/paper_skeleton.md ; measured numbers: Docs/paper_figures.md
%
% DUAL USE:
%   * Standalone: compiles as-is on Overleaf (the wrapper below supplies the
%     preamble, \begin{document}, and a manual bibliography).
%   * \input into a parent manuscript: define \theoryhardwareEMBED before the
%     \input and the standalone wrapper + bibliography are skipped, e.g.
%       \newcommand{\theoryhardwareEMBED}{}
%       % Vmem-phi -- Full draft sections: Corruption suite, Corruption algorithms, Theory, Hardware
% Companion outline: Docs/paper_skeleton.md ; measured numbers: Docs/paper_figures.md
%
% DUAL USE:
%   * Standalone: compiles as-is on Overleaf (the wrapper below supplies the
%     preamble, \begin{document}, and a manual bibliography).
%   * \input into a parent manuscript: define \theoryhardwareEMBED before the
%     \input and the standalone wrapper + bibliography are skipped, e.g.
%       \newcommand{\theoryhardwareEMBED}{}
%       \input{paper_sections_theory_hardware}
%
% Citation keys are placeholders; the embedded thebibliography below gives them
% resolvable (if approximate) entries so the standalone build is clean.

\ifdefined\theoryhardwareEMBED\else
\documentclass[11pt,a4paper]{article}
\usepackage[utf8]{inputenc}
\usepackage[margin=1in]{geometry}
\usepackage{amsmath}
\usepackage{amssymb}
\usepackage{booktabs}
\usepackage{tabularx}
\usepackage{microtype}
\begin{document}
\fi

\section{The Event-Camera Corruption Suite}
\label{sec:corruptions}

A benchmark for out-of-distribution (OOD) detection is only as meaningful as the
distribution shifts it tests against. We therefore define a controlled, reproducible
battery of six corruptions that model the dominant failure modes of event-camera
front-ends, each applied at five monotonically increasing severities. Together with the
clean baseline this yields $6\times5+1=31$ runs per detector. The suite is designed so that
each corruption perturbs a \emph{structurally distinct} axis of the input---DC offset,
localized burst, temporal order, global rate, polarity, or spatial support---which is what
makes the closed-form analysis of Section~\ref{sec:theory} possible: a detector's behaviour
on each corruption can be predicted from which axis it moves.

\subsection{Input representation}
\label{sec:corruptions-input}

Every corruption operates directly on the model's input tensor rather than on the raw event
stream, so that it composes with any downstream architecture and is label-free (no detection
annotations are consulted). A sequence is a stack of $N$ event-histogram frames of shape
$(20, H, W)$, with $H\times W = 240\times304$ for Prophesee Gen1. The $20$ channels encode
two polarities $\times$ ten time bins of width $50\,\mu\text{s}$: channels $0$--$9$ are the
ON-polarity bins (oldest to newest) and channels $10$--$19$ the OFF-polarity bins. Counts
are stored as \texttt{uint8}; corruptions that can increase counts promote to a wider integer
or float type internally and clip back to $[0,255]$, so no silent wrap-around occurs.

\subsection{Design principles}
\label{sec:corruptions-design}

Three properties hold across the suite. \textbf{(i) Reproducibility:} every random decision
(pixel coordinates, burst centres, flip masks, shift amounts, rate direction) is drawn from a
single seeded NumPy generator, so a (corruption, severity, seed) triple is deterministic.
\textbf{(ii) Hardware agnosticism:} each operation is expressed through an array module that
dispatches to NumPy or CuPy, so the identical corruption runs on CPU or GPU.
\textbf{(iii) Severity monotonicity:} severity $1$ is a mild perturbation and severity $5$ a
severe one, with parameters chosen to increase monotonically (Section~\ref{sec:corr-algos}),
enabling the Spearman-$\rho$ monotonicity analysis of the results section. Table~\ref{tab:corr-taxonomy}
summarizes the failure mode and perturbed axis for each corruption.

\begin{table}[t]
  \centering
  \small
  \begin{tabularx}{\linewidth}{@{}l l X@{}}
    \toprule
    Corruption & Physical failure mode & Axis of the input perturbed \\
    \midrule
    \texttt{hot\_pixel}        & defective / stuck pixels        & constant DC offset at fixed pixel locations \\
    \texttt{event\_flood}      & EMI / sensor saturation         & localized count bursts in space and time \\
    \texttt{temporal\_jitter}  & clock noise / timestamp error   & time-bin ordering (per polarity) \\
    \texttt{event\_rate\_shift}& event-density domain gap         & global multiplicative rate \\
    \texttt{polarity\_flip}    & ON/OFF misclassification         & polarity sign on a fraction of frames \\
    \texttt{spatial\_dropout}  & dead-pixel clusters / occlusion & spatial support (rectangular regions zeroed) \\
    \bottomrule
  \end{tabularx}
  \caption{The six event-camera corruptions, the sensor failure each models, and the
  structural axis of the input histogram it perturbs. Each axis maps to a distinct,
  predictable signature in the membrane moments (Section~\ref{sec:theory}).}
  \label{tab:corr-taxonomy}
\end{table}

\section{Corruption Algorithms}
\label{sec:corr-algos}

We now specify each corruption and its severity schedule. All operate on a histogram
$X \in \{0,\dots,255\}^{N\times 20\times H\times W}$ and return a corrupted copy of the same
shape and dtype.

\paragraph{Hot pixel.}
A fixed set of $n$ pixel locations is sampled uniformly at random and a constant count $\delta$
is added to \emph{every} time bin of \emph{every} frame at those locations, modelling pixels
that fire continuously regardless of scene content. The accumulation is performed in
\texttt{int16} and clipped to $[0,255]$. Severity raises both the number of hot pixels and the
injected count: $(n,\delta) \in \{(10,20),(30,40),(80,80),(150,140),(300,200)\}$ for
severities $1$--$5$.

\paragraph{Event flood.}
$b$ burst centres are placed at random frames and random spatial patches; at each, a count
$\delta$ is added inside a patch covering a fraction $f$ of the array, over the frames within
$\pm 2$ of the centre. This models electromagnetic interference or saturation affecting a
localized region for a short time. Severity raises the number of bursts, the patch fraction,
and the injected count: $(b,f,\delta) \in \{(2,0.05,60),(4,0.08,100),(8,0.12,150),(15,0.18,200),(25,0.25,255)\}$.

\paragraph{Temporal jitter.}
The ON-polarity bin block (channels $0$--$9$) and the OFF block (channels $10$--$19$) are each
rolled along the time-bin axis by an independent random integer shift drawn from
$[-s_{\max}, s_{\max}]$, with vacated bins zero-filled. This models clock noise / timestamp
inaccuracy that misassigns events to neighbouring time bins; crucially it permutes time order
while leaving the per-channel marginal counts almost unchanged. The maximum shift grows with
severity: $s_{\max} \in \{1,2,3,5,8\}$ bins (one bin $=50\,\mu\text{s}$).

\paragraph{Event-rate shift.}
All counts are multiplied by a single factor and clipped, modelling a domain gap in event
density (low- vs.\ high-texture or slow vs.\ fast motion). A seeded coin flip selects an
under-rate or over-rate factor, so the corruption is two-sided. The factor pair
$(\text{under},\text{over})$ widens with severity:
$\{(0.80,1.20),(0.65,1.40),(0.50,1.65),(0.35,2.00),(0.20,2.50)\}$.

\paragraph{Polarity flip.}
For a random fraction $p$ of frames, the ON and OFF channel blocks are swapped, modelling
sensor misclassification of brightness-change direction. The fraction of affected frames grows
with severity: $p \in \{0.05,0.10,0.20,0.35,0.50\}$.

\paragraph{Spatial dropout.}
$k$ rectangular regions, each covering fractions $(f_h, f_w)$ of the height and width, are
zeroed across all channels and all frames, modelling dead-pixel clusters or persistent
structured occlusion. Severity raises the number and size of regions:
$(k, f_h, f_w) \in \{(1,0.05,0.05),(1,0.10,0.10),(2,0.10,0.10),(2,0.15,0.15),(3,0.20,0.20)\}$.

\begin{table}[t]
  \centering
  \small
  \begin{tabularx}{\linewidth}{@{}l X X X X X@{}}
    \toprule
    Corruption & Sev.\ 1 & Sev.\ 2 & Sev.\ 3 & Sev.\ 4 & Sev.\ 5 \\
    \midrule
    \texttt{hot\_pixel} $(n,\delta)$        & $10,20$  & $30,40$  & $80,80$   & $150,140$ & $300,200$ \\
    \texttt{event\_flood} $(b,f,\delta)$    & $2,.05,60$ & $4,.08,100$ & $8,.12,150$ & $15,.18,200$ & $25,.25,255$ \\
    \texttt{temporal\_jitter} $s_{\max}$    & $1$ & $2$ & $3$ & $5$ & $8$ \\
    \texttt{event\_rate\_shift} (u, o)      & $.80,1.20$ & $.65,1.40$ & $.50,1.65$ & $.35,2.00$ & $.20,2.50$ \\
    \texttt{polarity\_flip} $p$             & $.05$ & $.10$ & $.20$ & $.35$ & $.50$ \\
    \texttt{spatial\_dropout} $(k,f_h,f_w)$ & $1,.05,.05$ & $1,.10,.10$ & $2,.10,.10$ & $2,.15,.15$ & $3,.20,.20$ \\
    \bottomrule
  \end{tabularx}
  \caption{Severity schedules for the six corruptions. Parameters increase monotonically with
  severity. Symbols: $n$ hot pixels, $\delta$ injected count, $b$ burst centres, $f$ patch
  fraction, $s_{\max}$ max bin shift, $(u,o)$ under/over rate factors, $p$ flipped-frame
  fraction, $k$ dropout regions, $(f_h,f_w)$ region height/width fractions.}
  \label{tab:corr-severity}
\end{table}

\section{A Leaky-Integrator Theory of Membrane-Statistics OOD Detection}
\label{sec:theory}

A benchmark that reports which corruptions a detector catches, but not \emph{why},
invites the obvious reviewer objection: the method may be exploiting dataset-specific
artifacts rather than a property of the underlying neuron model. This section removes that
objection. We show that the detectability of each corruption is not an empirical accident
but a closed-form consequence of the Parametric Leaky Integrate-and-Fire (PLIF) state
equation. The membrane potential is a \emph{leaky exponential accumulator} of the input
current; each corruption perturbs that current in a structurally distinct way, and the
perturbation propagates to a predictable change in the first three moments of the
sub-threshold trace. The resulting predictions match the measured per-corruption AUROC
ordering, turning the benchmark into a principled framework.

\subsection{The PLIF state equation and its closed-form solution}
\label{sec:theory-state}

The PLIF neuron, as instantiated by SpikingJelly's \texttt{MultiStepParametricLIFNode}
\cite{fang2021plif}, maintains a sub-threshold membrane potential $V$ that integrates an
input current with a learnable leak. In discrete multi-step form, the update for a single
channel is
%
\begin{equation}
  V[t] \;=\; \Big(1 - \tfrac{1}{\tau}\Big)\, V[t-1] \;+\; W\, X[t],
  \label{eq:plif}
\end{equation}
%
where $\tau > 1$ is the membrane time constant (learned per layer), $W$ is the synaptic
weight applied to the input event-histogram current $X[t]$ at time step $t$, and a spike is
emitted with a reset whenever $V$ crosses the firing threshold $\theta$. Writing the leak
factor as $\lambda \equiv 1 - 1/\tau \in (0,1)$ and unrolling \eqref{eq:plif} from a
quiescent initial state $V[-1]=0$ over the $T$ SNN time steps, the sub-threshold trace is an
exponentially weighted history of the input:
%
\begin{equation}
  V[t] \;=\; \sum_{k=0}^{t} \lambda^{\,k}\, W\, X[t-k].
  \label{eq:unroll}
\end{equation}
%
Equation~\eqref{eq:unroll} is the object the entire method rests on: $V$ is a \emph{causal,
leaky linear filter} of the event stream, with an effective memory horizon of $\sim\tau$
steps. Two immediate consequences frame everything that follows. First, $V$ is \emph{not} a
memoryless function of the current frame---it carries the temporal structure of recent
input, which is why temporally-structured corruptions leave a signature even when the
marginal input distribution is unchanged. Second, because the map from $X$ to $V$ is linear,
the \emph{moments} of $V$ inherit closed-form dependence on the moments of $X$, which we now
make explicit.

\subsection{The feature: per-channel membrane moments}
\label{sec:theory-feature}

For each channel $c$ in each of the four PLIF layers, we apply Global Average Pooling (GAP)
over the spatial dimensions of $V$ and compute the first three central moments of the pooled
trace across the $T$ steps:
%
\begin{equation}
  \varphi_c \;=\; \big[\, \mu_c,\; \sigma_c^2,\; \kappa_c \,\big],
  \qquad
  \mu_c = \mathbb{E}[V_c],\quad
  \sigma_c^2 = \mathrm{Var}[V_c],\quad
  \kappa_c = \frac{\mathbb{E}\big[(V_c-\mu_c)^4\big]}{\sigma_c^4}.
  \label{eq:phi}
\end{equation}
%
Stacking over $704$ channels across the four layers yields the $2112$-dimensional descriptor
$\varphi$. The mean captures sustained drive, the variance captures fluctuation energy, and
the kurtosis captures burstiness (heavy-tailed, intermittent excitation). We show next that
these three coordinates are individually informative about \emph{different} corruption
families.

\subsection{Steady-state moments under a stationary input}
\label{sec:theory-moments}

Suppose the per-channel input current is wide-sense stationary with mean $m=\mathbb{E}[X]$,
variance $s^2=\mathrm{Var}[X]$, and autocovariance $\gamma(j)=\mathrm{Cov}(X[t],X[t-j])$.
Taking expectations of the geometric sum \eqref{eq:unroll} in the long-horizon limit
($\sum_k \lambda^k \to 1/(1-\lambda)$) gives the steady-state membrane \emph{mean}
%
\begin{equation}
  \mathbb{E}[V] \;=\; \frac{W m}{1-\lambda} \;=\; \tau\, W m,
  \label{eq:mean}
\end{equation}
%
so the mean is the input rate \emph{amplified by the time constant} $\tau$. The steady-state
\emph{variance}, assuming weak temporal correlation ($\gamma(j)\!\approx\!0$ for $j>0$), is
%
\begin{equation}
  \mathrm{Var}[V] \;=\; (W s)^2 \sum_{k=0}^{\infty} \lambda^{2k}
  \;=\; \frac{(W s)^2}{1-\lambda^2}
  \;=\; \frac{\tau^2}{2\tau-1}\,(W s)^2,
  \label{eq:var}
\end{equation}
%
i.e.\ the membrane variance tracks the \emph{input variance}, scaled by a leak-dependent
constant. Retaining the autocovariance term generalizes \eqref{eq:var} to
%
\begin{equation}
  \mathrm{Var}[V] \;=\; \frac{(W s)^2}{1-\lambda^2}
  \;+\; 2 W^2 \sum_{j=1}^{\infty} \gamma(j)\,\frac{\lambda^{\,j}}{1-\lambda^2},
  \label{eq:var-auto}
\end{equation}
%
which is the key to temporal corruptions: the second term depends on the \emph{time
ordering} of the input, not its marginal distribution. A corruption that preserves $m$ and
$s^2$ but scrambles $\gamma(j)$ moves $\mathrm{Var}[V]$ (and, through higher-order analogues,
the kurtosis and the autocorrelation of $V(t)$) while leaving the naive first-moment picture
untouched. Equations~\eqref{eq:mean}--\eqref{eq:var-auto} are the three levers every
corruption pulls; the rest of the section reads off, corruption by corruption, which lever
moves and therefore which feature detects it.

\subsection{Per-corruption analysis}
\label{sec:theory-percorruption}

\paragraph{Hot-pixel --- DC saturation (predicted AUROC $\approx 1$).}
A stuck/hot pixel injects a persistent constant $\delta$ into $X[t]$ at fixed spatial
locations for the affected channels. Adding a constant $\delta$ to the input in
\eqref{eq:mean} shifts the membrane mean by $\tau W \delta$, an unbounded accumulation that
drives $V$ monotonically toward---and clamps it at---the firing threshold $\theta$. The
affected channels' $\mu_c$ therefore saturate at a value disjoint from the clean range,
producing two non-overlapping clusters in $\varphi$-space. Detection is not merely easy but
\emph{mathematically inevitable}, and it is monotone in severity because larger $\delta$
saturates more channels. This is exactly the measured behaviour (perfect AUROC and Spearman
monotonicity).

\paragraph{Temporal jitter --- autocorrelation destruction (temporal/deep features only).}
Temporal jitter permutes the time order of $X[t]$ within a window. A permutation is
measure-preserving: it leaves the marginal mean $m$ and variance $s^2$---and hence
\eqref{eq:mean} and the \emph{first} term of \eqref{eq:var-auto}---unchanged. What it
destroys is the autocovariance $\gamma(j)$, i.e.\ the second term of \eqref{eq:var-auto} and,
more sharply, the lag-1 autocorrelation of the membrane trace $V(t)$ itself. The leaky filter
\eqref{eq:unroll} is precisely an autocorrelation-generating operation, so jitter is visible
in the \emph{temporal} structure of $V$ even though it is invisible to the static marginal
mean. This predicts, and the data confirm, that jitter is caught by temporal statistics
(lag-1 autocorrelation) and by deeper layers whose larger effective $\tau$ integrate
longer-range temporal structure---the layer-4 slice is the strongest single view for jitter.

\paragraph{Event-rate shift --- global affine drive (single activity scalar).}
An event-rate shift multiplies the global input rate $m \to \alpha m$. By \eqref{eq:mean}
this scales the membrane mean of \emph{every} channel by the same factor $\alpha$, an
approximately affine displacement of $\varphi$ along a single global-activity direction. The
corruption is thus detectable by a one-dimensional projection---a global activity
scalar---which is why a single summary statistic recovers most of the signal, and why the
per-frame AUROC shows a phase transition once $\alpha$ leaves the clean operating band.

\paragraph{Spatial dropout --- variance loss hidden by GAP (anti-detectable under one-sided
scoring).}
Spatial dropout zeros a contiguous spatial \emph{region} of $X[t]$. Within an affected
channel this lowers the spatial input variance and, through \eqref{eq:var}, the membrane
variance $\sigma_c^2$, while leaving the channel mean $\mu_c$ comparatively intact (the
surviving region still drives the neuron). Variance is therefore the single most informative
moment for this corruption---but GAP pools away the \emph{spatial} dimension before the
moments are taken, so the dominant signature (which spatial regions went silent) is
discarded, and what remains makes the membrane look \emph{quieter and more regular}, i.e.\
closer to the clean mean. Under one-sided distance-from-clean scoring this inverts the
detector (AUROC $< 0.5$): the corruption is anti-detectable not because the information is
absent from $V$ but because GAP plus a one-sided rule throws it away. This is the formal
motivation for the spatial-dispersion features (\texttt{phi\_spatial}) and for two-sided
scoring (Section~\ref{sec:method}).

\paragraph{Event flood --- uniform mean shift indistinguishable from a busy scene
(near-chance per frame).}
Event flood adds approximately uniform noise across all pixels, raising $m$ uniformly. By
\eqref{eq:mean} this lifts the mean of \emph{all} channels by nearly the same amount---the
identical signature a genuinely busy, high-activity clean scene produces. Because the clean
distribution already contains such high-activity scenes, the flooded frames are not separable
from the clean tail on the marginal moments alone, giving near-chance per-frame AUROC. The
information that \emph{does} distinguish flood---that the elevated activity is spatially
unstructured (it fills space rather than tracking objects)---again lives in the spatial
dispersion GAP removes, and in the temporal \emph{consistency} of the bias across a sequence.
This predicts the two recoveries we observe: spatial features, and sequence-level
aggregation, which lifts flood to near-perfect AUROC once the per-frame bias is integrated
over a whole recording.

\paragraph{Polarity flip --- invariance by construction ($\approx 0.5$ ceiling, a
membrane-side dead end).}
Polarity flip inverts the ON/OFF sign of the events. The upstream detector was trained to be
approximately polarity-symmetric; to the extent that symmetry holds, $W X$ is near-invariant
to the flip, so $V$---and therefore $\varphi$---barely moves, and the AUROC ceiling is
$\approx 0.5$ \emph{by construction}, not through any deficiency of the detector. The small,
consistent residual below chance (a few points) reflects imperfect symmetry, i.e.\ a faint
directional shift rather than noise. The honest conclusion is that polarity\_flip is largely
undetectable from the membrane and its signal lives on the \emph{input} side (the
positive/negative event-count balance the flip inverts); we scope it accordingly rather than
chase it.

\subsection{Why GAP causes the inversions, and what the theory prescribes}
\label{sec:theory-taxonomy}

The three hard corruptions---spatial\_dropout, event\_flood, polarity\_flip---share a single
mechanism in the analysis above: their discriminative signature lives in structure that
GAP-then-moments discards (spatial dispersion for the first two, input-side sign balance for
the third), while the surviving marginal moments are preserved or even pushed toward the
clean mean. This yields a clean, falsifiable taxonomy that organizes the entire results
section:
%
\begin{itemize}
  \item \textbf{Magnitude corruptions} (hot\_pixel, event\_rate\_shift, event\_flood-as-mean,
        temporal\_jitter via variance) move the membrane moments \emph{away} from clean and
        are caught by distance-based detectors on $\varphi$.
  \item \textbf{Structure corruptions} (spatial\_dropout, and the spatial component of
        event\_flood) leave the marginal moments intact and require either the
        spatial-dispersion features GAP discards or a scoring rule that flags ``too
        typical,'' not only ``too far.''
  \item \textbf{Invariant corruptions} (polarity\_flip) are near-orthogonal to the membrane
        and are an input-side, not a membrane-side, detection problem.
\end{itemize}
%
The theory therefore does three things at once: it \emph{explains} the easy wins as
inevitable, it \emph{predicts} the inversions from first principles, and it \emph{prescribes}
the architecture that addresses them---the spatial branch, two-sided scoring, and sequence
aggregation that the MDD detector (Section~\ref{sec:method}) combines. Crucially, it does so
without any free parameters fit to the corruption labels: every prediction follows from the
PLIF state equation and the definition of the feature.

\section{The Neuromorphic Hardware Advantage}
\label{sec:hardware}

The empirical case for membrane-statistics OOD detection is incomplete without its central
positioning claim, which is not about accuracy at all. Existing OOD detectors are designed
for artificial neural networks executed on von Neumann hardware, where logits, gradients, and
repeated forward passes are cheap. Spiking networks are deployed for the opposite reason---to
run within the severe energy and compute budget of neuromorphic silicon---and in that regime
the standard OOD toolkit is not merely expensive but \emph{architecturally infeasible}.
Membrane statistics are the rare signal that is free precisely where every alternative is
impossible. This reframes the contribution from ``a competitive OOD score'' to ``the only OOD
method that runs on the hardware the model targets.''

\subsection{The membrane potential is a native, already-maintained register}
\label{sec:hardware-register}

On the major neuromorphic platforms---Intel Loihi \cite{davies2018loihi}, the BrainScaleS
analog system \cite{schemmel2010brainscales}, and SpiNNaker \cite{furber2014spinnaker}---each
neuron's membrane potential $V$ is \emph{physical state the substrate already maintains}. It
is not an intermediate that must be recomputed for our purposes; it is the quantity the neuron
updates on every time step in order to decide whether to spike, held in a local register
(digital cores) or a capacitor voltage (analog). Producing our descriptor $\varphi$ from it
requires only a read of that state and a reduction to three per-channel moments---an
$O(\text{channels})$ operation with no additional synaptic integration. In the cost accounting
that matters on-chip---multiply--accumulate operations and memory traffic for weights---%
$\varphi$ adds \emph{zero MACs and zero weight fetches}: it is a byproduct of the forward pass
that has to happen anyway to emit spikes. This is the precise sense in which the method is
``zero additional compute,'' and it is a property of \emph{where} the signal comes from, not
an optimization we apply afterward.

\subsection{Why competing OOD methods cannot run on the target hardware}
\label{sec:hardware-table}

The dominant OOD families each presuppose a capability that the neuromorphic execution model
does not provide. We enumerate them not to claim our accuracy is higher, but to establish that
the comparison cannot even take place on-chip.

\begin{table}[t]
  \centering
  \small
  \begin{tabularx}{\linewidth}{@{}l X X@{}}
    \toprule
    Method & Required extra computation & Why infeasible on neuromorphic HW \\
    \midrule
    \textbf{Vmem-$\varphi$ (ours)} & native $V$ register read + moments & --- (it is free) \\
    MSP \cite{hendrycks2017msp} / Energy \cite{liu2020energy}
      & softmax / log-sum-exp over logits
      & cores emit spikes, not float logits; no on-chip softmax \\
    ODIN \cite{liang2018odin}
      & temperature softmax + input perturbation
      & perturbation needs an input-gradient pass the chip cannot take \\
    ReAct \cite{sun2021react}
      & activation clipping + second forward pass
      & re-forwarding violates the SNN energy budget \\
    ViM \cite{wang2022vim}
      & SVD / dense linear algebra on features
      & no dense float linear-algebra unit; data movement alone too costly \\
    GradNorm \cite{huang2021gradnorm}
      & gradient backpropagation at inference
      & neuromorphic inference has no backward pass \\
    kNN / Mahalanobis on \emph{ANN} features \cite{lee2018mahalanobis,sun2022knn}
      & an auxiliary ANN forward pass
      & requires a second, non-spiking network---defeating the SNN premise \\
    \bottomrule
  \end{tabularx}
  \caption{Standard OOD detectors each require a capability absent from the spiking,
  feed-forward, integer-event execution model; Vmem-$\varphi$ requires none because the signal
  is the neuron state itself.}
  \label{tab:hardware}
\end{table}

The pattern is uniform: every alternative needs either a floating-point readout layer
(softmax, SVD), a second pass (ReAct, ANN-feature methods), or a gradient (ODIN,
GradNorm)---none of which exist in the spiking, feed-forward, integer-event execution model.
Vmem-$\varphi$ needs none of them because the signal is the neuron state itself.

\subsection{The reframed contribution}
\label{sec:hardware-claim}

It follows that the benchmark of Section~\ref{sec:experiments} should be read not as
``$\varphi$ narrowly beats ANN OOD baselines on a GPU'' but as ``$\varphi$ provides usable OOD
detection \emph{inside the compute envelope of the deployment hardware}, where the listed
alternatives cannot be evaluated at all.'' This is a categorical claim about feasibility, and
it is what makes the work a distinct contribution rather than another row in an OOD
leaderboard: for spiking networks on neuromorphic silicon, membrane-statistics OOD detection
is not the best option among many---under a strict on-chip compute budget it is, to our
knowledge, the only option that is free.

\subsection{Scope and honesty caveat}
\label{sec:hardware-caveat}

We measure $\varphi$ in simulation (SpikingJelly on a GPU), not on neuromorphic silicon. Our
cost claim is therefore \emph{architectural}---it follows from $V$ being a register the
substrate already maintains and from $\varphi$ adding no synaptic operations---rather than a
wall-clock or joule measurement on-chip. We state this explicitly: the contribution is that
the signal is free \emph{by construction} on the target hardware, and an on-silicon energy
measurement is identified as future work. We likewise restrict the infeasibility argument to
the standard execution models of current neuromorphic platforms; co-processor or hybrid
configurations could in principle host some of the listed methods at a cost, which only
sharpens the comparison rather than changing its direction.

% --- Standalone bibliography (skipped when embedded). Manual thebibliography so
%     the standalone build needs no .bib file and no bibtex pass. Replace with a
%     real \bibliography{...} when merging into the full manuscript.
\ifdefined\theoryhardwareEMBED\else
\begin{thebibliography}{99}

\bibitem{fang2021plif}
W.~Fang, Z.~Yu, Y.~Chen, T.~Masquelier, T.~Huang, and Y.~Tian.
\newblock Incorporating learnable membrane time constant to enhance learning of
  spiking neural networks.
\newblock In \emph{ICCV}, 2021.

\bibitem{davies2018loihi}
M.~Davies et~al.
\newblock Loihi: A neuromorphic manycore processor with on-chip learning.
\newblock \emph{IEEE Micro}, 38(1):82--99, 2018.

\bibitem{schemmel2010brainscales}
J.~Schemmel, D.~Br\"uderle, A.~Gr\"ubl, M.~Hock, K.~Meier, and S.~Millner.
\newblock A wafer-scale neuromorphic hardware system for large-scale neural
  modeling.
\newblock In \emph{ISCAS}, pages 1947--1950, 2010.

\bibitem{furber2014spinnaker}
S.~B. Furber, F.~Galluppi, S.~Temple, and L.~A. Plana.
\newblock The SpiNNaker project.
\newblock \emph{Proceedings of the IEEE}, 102(5):652--665, 2014.

\bibitem{hendrycks2017msp}
D.~Hendrycks and K.~Gimpel.
\newblock A baseline for detecting misclassified and out-of-distribution
  examples in neural networks.
\newblock In \emph{ICLR}, 2017.

\bibitem{liu2020energy}
W.~Liu, X.~Wang, J.~D. Owens, and Y.~Li.
\newblock Energy-based out-of-distribution detection.
\newblock In \emph{NeurIPS}, 2020.

\bibitem{liang2018odin}
S.~Liang, Y.~Li, and R.~Srikant.
\newblock Enhancing the reliability of out-of-distribution image detection in
  neural networks.
\newblock In \emph{ICLR}, 2018.

\bibitem{sun2021react}
Y.~Sun, C.~Guo, and Y.~Li.
\newblock ReAct: Out-of-distribution detection with rectified activations.
\newblock In \emph{NeurIPS}, 2021.

\bibitem{wang2022vim}
H.~Wang, Z.~Li, L.~Feng, and W.~Zhang.
\newblock ViM: Out-of-distribution with virtual-logit matching.
\newblock In \emph{CVPR}, 2022.

\bibitem{huang2021gradnorm}
R.~Huang, A.~Geng, and Y.~Li.
\newblock On the importance of gradients for detecting distributional shifts in
  the wild.
\newblock In \emph{NeurIPS}, 2021.

\bibitem{lee2018mahalanobis}
K.~Lee, K.~Lee, H.~Lee, and J.~Shin.
\newblock A simple unified framework for detecting out-of-distribution samples
  and adversarial attacks.
\newblock In \emph{NeurIPS}, 2018.

\bibitem{sun2022knn}
Y.~Sun, Y.~Ming, X.~Zhu, and Y.~Li.
\newblock Out-of-distribution detection with deep nearest neighbors.
\newblock In \emph{ICML}, 2022.

\end{thebibliography}
\end{document}
\fi

%
% Citation keys are placeholders; the embedded thebibliography below gives them
% resolvable (if approximate) entries so the standalone build is clean.

\ifdefined\theoryhardwareEMBED\else
\documentclass[11pt,a4paper]{article}
\usepackage[utf8]{inputenc}
\usepackage[margin=1in]{geometry}
\usepackage{amsmath}
\usepackage{amssymb}
\usepackage{booktabs}
\usepackage{tabularx}
\usepackage{microtype}
\begin{document}
\fi

\section{The Event-Camera Corruption Suite}
\label{sec:corruptions}

A benchmark for out-of-distribution (OOD) detection is only as meaningful as the
distribution shifts it tests against. We therefore define a controlled, reproducible
battery of six corruptions that model the dominant failure modes of event-camera
front-ends, each applied at five monotonically increasing severities. Together with the
clean baseline this yields $6\times5+1=31$ runs per detector. The suite is designed so that
each corruption perturbs a \emph{structurally distinct} axis of the input---DC offset,
localized burst, temporal order, global rate, polarity, or spatial support---which is what
makes the closed-form analysis of Section~\ref{sec:theory} possible: a detector's behaviour
on each corruption can be predicted from which axis it moves.

\subsection{Input representation}
\label{sec:corruptions-input}

Every corruption operates directly on the model's input tensor rather than on the raw event
stream, so that it composes with any downstream architecture and is label-free (no detection
annotations are consulted). A sequence is a stack of $N$ event-histogram frames of shape
$(20, H, W)$, with $H\times W = 240\times304$ for Prophesee Gen1. The $20$ channels encode
two polarities $\times$ ten time bins of width $50\,\mu\text{s}$: channels $0$--$9$ are the
ON-polarity bins (oldest to newest) and channels $10$--$19$ the OFF-polarity bins. Counts
are stored as \texttt{uint8}; corruptions that can increase counts promote to a wider integer
or float type internally and clip back to $[0,255]$, so no silent wrap-around occurs.

\subsection{Design principles}
\label{sec:corruptions-design}

Three properties hold across the suite. \textbf{(i) Reproducibility:} every random decision
(pixel coordinates, burst centres, flip masks, shift amounts, rate direction) is drawn from a
single seeded NumPy generator, so a (corruption, severity, seed) triple is deterministic.
\textbf{(ii) Hardware agnosticism:} each operation is expressed through an array module that
dispatches to NumPy or CuPy, so the identical corruption runs on CPU or GPU.
\textbf{(iii) Severity monotonicity:} severity $1$ is a mild perturbation and severity $5$ a
severe one, with parameters chosen to increase monotonically (Section~\ref{sec:corr-algos}),
enabling the Spearman-$\rho$ monotonicity analysis of the results section. Table~\ref{tab:corr-taxonomy}
summarizes the failure mode and perturbed axis for each corruption.

\begin{table}[t]
  \centering
  \small
  \begin{tabularx}{\linewidth}{@{}l l X@{}}
    \toprule
    Corruption & Physical failure mode & Axis of the input perturbed \\
    \midrule
    \texttt{hot\_pixel}        & defective / stuck pixels        & constant DC offset at fixed pixel locations \\
    \texttt{event\_flood}      & EMI / sensor saturation         & localized count bursts in space and time \\
    \texttt{temporal\_jitter}  & clock noise / timestamp error   & time-bin ordering (per polarity) \\
    \texttt{event\_rate\_shift}& event-density domain gap         & global multiplicative rate \\
    \texttt{polarity\_flip}    & ON/OFF misclassification         & polarity sign on a fraction of frames \\
    \texttt{spatial\_dropout}  & dead-pixel clusters / occlusion & spatial support (rectangular regions zeroed) \\
    \bottomrule
  \end{tabularx}
  \caption{The six event-camera corruptions, the sensor failure each models, and the
  structural axis of the input histogram it perturbs. Each axis maps to a distinct,
  predictable signature in the membrane moments (Section~\ref{sec:theory}).}
  \label{tab:corr-taxonomy}
\end{table}

\section{Corruption Algorithms}
\label{sec:corr-algos}

We now specify each corruption and its severity schedule. All operate on a histogram
$X \in \{0,\dots,255\}^{N\times 20\times H\times W}$ and return a corrupted copy of the same
shape and dtype.

\paragraph{Hot pixel.}
A fixed set of $n$ pixel locations is sampled uniformly at random and a constant count $\delta$
is added to \emph{every} time bin of \emph{every} frame at those locations, modelling pixels
that fire continuously regardless of scene content. The accumulation is performed in
\texttt{int16} and clipped to $[0,255]$. Severity raises both the number of hot pixels and the
injected count: $(n,\delta) \in \{(10,20),(30,40),(80,80),(150,140),(300,200)\}$ for
severities $1$--$5$.

\paragraph{Event flood.}
$b$ burst centres are placed at random frames and random spatial patches; at each, a count
$\delta$ is added inside a patch covering a fraction $f$ of the array, over the frames within
$\pm 2$ of the centre. This models electromagnetic interference or saturation affecting a
localized region for a short time. Severity raises the number of bursts, the patch fraction,
and the injected count: $(b,f,\delta) \in \{(2,0.05,60),(4,0.08,100),(8,0.12,150),(15,0.18,200),(25,0.25,255)\}$.

\paragraph{Temporal jitter.}
The ON-polarity bin block (channels $0$--$9$) and the OFF block (channels $10$--$19$) are each
rolled along the time-bin axis by an independent random integer shift drawn from
$[-s_{\max}, s_{\max}]$, with vacated bins zero-filled. This models clock noise / timestamp
inaccuracy that misassigns events to neighbouring time bins; crucially it permutes time order
while leaving the per-channel marginal counts almost unchanged. The maximum shift grows with
severity: $s_{\max} \in \{1,2,3,5,8\}$ bins (one bin $=50\,\mu\text{s}$).

\paragraph{Event-rate shift.}
All counts are multiplied by a single factor and clipped, modelling a domain gap in event
density (low- vs.\ high-texture or slow vs.\ fast motion). A seeded coin flip selects an
under-rate or over-rate factor, so the corruption is two-sided. The factor pair
$(\text{under},\text{over})$ widens with severity:
$\{(0.80,1.20),(0.65,1.40),(0.50,1.65),(0.35,2.00),(0.20,2.50)\}$.

\paragraph{Polarity flip.}
For a random fraction $p$ of frames, the ON and OFF channel blocks are swapped, modelling
sensor misclassification of brightness-change direction. The fraction of affected frames grows
with severity: $p \in \{0.05,0.10,0.20,0.35,0.50\}$.

\paragraph{Spatial dropout.}
$k$ rectangular regions, each covering fractions $(f_h, f_w)$ of the height and width, are
zeroed across all channels and all frames, modelling dead-pixel clusters or persistent
structured occlusion. Severity raises the number and size of regions:
$(k, f_h, f_w) \in \{(1,0.05,0.05),(1,0.10,0.10),(2,0.10,0.10),(2,0.15,0.15),(3,0.20,0.20)\}$.

\begin{table}[t]
  \centering
  \small
  \begin{tabularx}{\linewidth}{@{}l X X X X X@{}}
    \toprule
    Corruption & Sev.\ 1 & Sev.\ 2 & Sev.\ 3 & Sev.\ 4 & Sev.\ 5 \\
    \midrule
    \texttt{hot\_pixel} $(n,\delta)$        & $10,20$  & $30,40$  & $80,80$   & $150,140$ & $300,200$ \\
    \texttt{event\_flood} $(b,f,\delta)$    & $2,.05,60$ & $4,.08,100$ & $8,.12,150$ & $15,.18,200$ & $25,.25,255$ \\
    \texttt{temporal\_jitter} $s_{\max}$    & $1$ & $2$ & $3$ & $5$ & $8$ \\
    \texttt{event\_rate\_shift} (u, o)      & $.80,1.20$ & $.65,1.40$ & $.50,1.65$ & $.35,2.00$ & $.20,2.50$ \\
    \texttt{polarity\_flip} $p$             & $.05$ & $.10$ & $.20$ & $.35$ & $.50$ \\
    \texttt{spatial\_dropout} $(k,f_h,f_w)$ & $1,.05,.05$ & $1,.10,.10$ & $2,.10,.10$ & $2,.15,.15$ & $3,.20,.20$ \\
    \bottomrule
  \end{tabularx}
  \caption{Severity schedules for the six corruptions. Parameters increase monotonically with
  severity. Symbols: $n$ hot pixels, $\delta$ injected count, $b$ burst centres, $f$ patch
  fraction, $s_{\max}$ max bin shift, $(u,o)$ under/over rate factors, $p$ flipped-frame
  fraction, $k$ dropout regions, $(f_h,f_w)$ region height/width fractions.}
  \label{tab:corr-severity}
\end{table}

\section{A Leaky-Integrator Theory of Membrane-Statistics OOD Detection}
\label{sec:theory}

A benchmark that reports which corruptions a detector catches, but not \emph{why},
invites the obvious reviewer objection: the method may be exploiting dataset-specific
artifacts rather than a property of the underlying neuron model. This section removes that
objection. We show that the detectability of each corruption is not an empirical accident
but a closed-form consequence of the Parametric Leaky Integrate-and-Fire (PLIF) state
equation. The membrane potential is a \emph{leaky exponential accumulator} of the input
current; each corruption perturbs that current in a structurally distinct way, and the
perturbation propagates to a predictable change in the first three moments of the
sub-threshold trace. The resulting predictions match the measured per-corruption AUROC
ordering, turning the benchmark into a principled framework.

\subsection{The PLIF state equation and its closed-form solution}
\label{sec:theory-state}

The PLIF neuron, as instantiated by SpikingJelly's \texttt{MultiStepParametricLIFNode}
\cite{fang2021plif}, maintains a sub-threshold membrane potential $V$ that integrates an
input current with a learnable leak. In discrete multi-step form, the update for a single
channel is
%
\begin{equation}
  V[t] \;=\; \Big(1 - \tfrac{1}{\tau}\Big)\, V[t-1] \;+\; W\, X[t],
  \label{eq:plif}
\end{equation}
%
where $\tau > 1$ is the membrane time constant (learned per layer), $W$ is the synaptic
weight applied to the input event-histogram current $X[t]$ at time step $t$, and a spike is
emitted with a reset whenever $V$ crosses the firing threshold $\theta$. Writing the leak
factor as $\lambda \equiv 1 - 1/\tau \in (0,1)$ and unrolling \eqref{eq:plif} from a
quiescent initial state $V[-1]=0$ over the $T$ SNN time steps, the sub-threshold trace is an
exponentially weighted history of the input:
%
\begin{equation}
  V[t] \;=\; \sum_{k=0}^{t} \lambda^{\,k}\, W\, X[t-k].
  \label{eq:unroll}
\end{equation}
%
Equation~\eqref{eq:unroll} is the object the entire method rests on: $V$ is a \emph{causal,
leaky linear filter} of the event stream, with an effective memory horizon of $\sim\tau$
steps. Two immediate consequences frame everything that follows. First, $V$ is \emph{not} a
memoryless function of the current frame---it carries the temporal structure of recent
input, which is why temporally-structured corruptions leave a signature even when the
marginal input distribution is unchanged. Second, because the map from $X$ to $V$ is linear,
the \emph{moments} of $V$ inherit closed-form dependence on the moments of $X$, which we now
make explicit.

\subsection{The feature: per-channel membrane moments}
\label{sec:theory-feature}

For each channel $c$ in each of the four PLIF layers, we apply Global Average Pooling (GAP)
over the spatial dimensions of $V$ and compute the first three central moments of the pooled
trace across the $T$ steps:
%
\begin{equation}
  \varphi_c \;=\; \big[\, \mu_c,\; \sigma_c^2,\; \kappa_c \,\big],
  \qquad
  \mu_c = \mathbb{E}[V_c],\quad
  \sigma_c^2 = \mathrm{Var}[V_c],\quad
  \kappa_c = \frac{\mathbb{E}\big[(V_c-\mu_c)^4\big]}{\sigma_c^4}.
  \label{eq:phi}
\end{equation}
%
Stacking over $704$ channels across the four layers yields the $2112$-dimensional descriptor
$\varphi$. The mean captures sustained drive, the variance captures fluctuation energy, and
the kurtosis captures burstiness (heavy-tailed, intermittent excitation). We show next that
these three coordinates are individually informative about \emph{different} corruption
families.

\subsection{Steady-state moments under a stationary input}
\label{sec:theory-moments}

Suppose the per-channel input current is wide-sense stationary with mean $m=\mathbb{E}[X]$,
variance $s^2=\mathrm{Var}[X]$, and autocovariance $\gamma(j)=\mathrm{Cov}(X[t],X[t-j])$.
Taking expectations of the geometric sum \eqref{eq:unroll} in the long-horizon limit
($\sum_k \lambda^k \to 1/(1-\lambda)$) gives the steady-state membrane \emph{mean}
%
\begin{equation}
  \mathbb{E}[V] \;=\; \frac{W m}{1-\lambda} \;=\; \tau\, W m,
  \label{eq:mean}
\end{equation}
%
so the mean is the input rate \emph{amplified by the time constant} $\tau$. The steady-state
\emph{variance}, assuming weak temporal correlation ($\gamma(j)\!\approx\!0$ for $j>0$), is
%
\begin{equation}
  \mathrm{Var}[V] \;=\; (W s)^2 \sum_{k=0}^{\infty} \lambda^{2k}
  \;=\; \frac{(W s)^2}{1-\lambda^2}
  \;=\; \frac{\tau^2}{2\tau-1}\,(W s)^2,
  \label{eq:var}
\end{equation}
%
i.e.\ the membrane variance tracks the \emph{input variance}, scaled by a leak-dependent
constant. Retaining the autocovariance term generalizes \eqref{eq:var} to
%
\begin{equation}
  \mathrm{Var}[V] \;=\; \frac{(W s)^2}{1-\lambda^2}
  \;+\; 2 W^2 \sum_{j=1}^{\infty} \gamma(j)\,\frac{\lambda^{\,j}}{1-\lambda^2},
  \label{eq:var-auto}
\end{equation}
%
which is the key to temporal corruptions: the second term depends on the \emph{time
ordering} of the input, not its marginal distribution. A corruption that preserves $m$ and
$s^2$ but scrambles $\gamma(j)$ moves $\mathrm{Var}[V]$ (and, through higher-order analogues,
the kurtosis and the autocorrelation of $V(t)$) while leaving the naive first-moment picture
untouched. Equations~\eqref{eq:mean}--\eqref{eq:var-auto} are the three levers every
corruption pulls; the rest of the section reads off, corruption by corruption, which lever
moves and therefore which feature detects it.

\subsection{Per-corruption analysis}
\label{sec:theory-percorruption}

\paragraph{Hot-pixel --- DC saturation (predicted AUROC $\approx 1$).}
A stuck/hot pixel injects a persistent constant $\delta$ into $X[t]$ at fixed spatial
locations for the affected channels. Adding a constant $\delta$ to the input in
\eqref{eq:mean} shifts the membrane mean by $\tau W \delta$, an unbounded accumulation that
drives $V$ monotonically toward---and clamps it at---the firing threshold $\theta$. The
affected channels' $\mu_c$ therefore saturate at a value disjoint from the clean range,
producing two non-overlapping clusters in $\varphi$-space. Detection is not merely easy but
\emph{mathematically inevitable}, and it is monotone in severity because larger $\delta$
saturates more channels. This is exactly the measured behaviour (perfect AUROC and Spearman
monotonicity).

\paragraph{Temporal jitter --- autocorrelation destruction (temporal/deep features only).}
Temporal jitter permutes the time order of $X[t]$ within a window. A permutation is
measure-preserving: it leaves the marginal mean $m$ and variance $s^2$---and hence
\eqref{eq:mean} and the \emph{first} term of \eqref{eq:var-auto}---unchanged. What it
destroys is the autocovariance $\gamma(j)$, i.e.\ the second term of \eqref{eq:var-auto} and,
more sharply, the lag-1 autocorrelation of the membrane trace $V(t)$ itself. The leaky filter
\eqref{eq:unroll} is precisely an autocorrelation-generating operation, so jitter is visible
in the \emph{temporal} structure of $V$ even though it is invisible to the static marginal
mean. This predicts, and the data confirm, that jitter is caught by temporal statistics
(lag-1 autocorrelation) and by deeper layers whose larger effective $\tau$ integrate
longer-range temporal structure---the layer-4 slice is the strongest single view for jitter.

\paragraph{Event-rate shift --- global affine drive (single activity scalar).}
An event-rate shift multiplies the global input rate $m \to \alpha m$. By \eqref{eq:mean}
this scales the membrane mean of \emph{every} channel by the same factor $\alpha$, an
approximately affine displacement of $\varphi$ along a single global-activity direction. The
corruption is thus detectable by a one-dimensional projection---a global activity
scalar---which is why a single summary statistic recovers most of the signal, and why the
per-frame AUROC shows a phase transition once $\alpha$ leaves the clean operating band.

\paragraph{Spatial dropout --- variance loss hidden by GAP (anti-detectable under one-sided
scoring).}
Spatial dropout zeros a contiguous spatial \emph{region} of $X[t]$. Within an affected
channel this lowers the spatial input variance and, through \eqref{eq:var}, the membrane
variance $\sigma_c^2$, while leaving the channel mean $\mu_c$ comparatively intact (the
surviving region still drives the neuron). Variance is therefore the single most informative
moment for this corruption---but GAP pools away the \emph{spatial} dimension before the
moments are taken, so the dominant signature (which spatial regions went silent) is
discarded, and what remains makes the membrane look \emph{quieter and more regular}, i.e.\
closer to the clean mean. Under one-sided distance-from-clean scoring this inverts the
detector (AUROC $< 0.5$): the corruption is anti-detectable not because the information is
absent from $V$ but because GAP plus a one-sided rule throws it away. This is the formal
motivation for the spatial-dispersion features (\texttt{phi\_spatial}) and for two-sided
scoring (Section~\ref{sec:method}).

\paragraph{Event flood --- uniform mean shift indistinguishable from a busy scene
(near-chance per frame).}
Event flood adds approximately uniform noise across all pixels, raising $m$ uniformly. By
\eqref{eq:mean} this lifts the mean of \emph{all} channels by nearly the same amount---the
identical signature a genuinely busy, high-activity clean scene produces. Because the clean
distribution already contains such high-activity scenes, the flooded frames are not separable
from the clean tail on the marginal moments alone, giving near-chance per-frame AUROC. The
information that \emph{does} distinguish flood---that the elevated activity is spatially
unstructured (it fills space rather than tracking objects)---again lives in the spatial
dispersion GAP removes, and in the temporal \emph{consistency} of the bias across a sequence.
This predicts the two recoveries we observe: spatial features, and sequence-level
aggregation, which lifts flood to near-perfect AUROC once the per-frame bias is integrated
over a whole recording.

\paragraph{Polarity flip --- invariance by construction ($\approx 0.5$ ceiling, a
membrane-side dead end).}
Polarity flip inverts the ON/OFF sign of the events. The upstream detector was trained to be
approximately polarity-symmetric; to the extent that symmetry holds, $W X$ is near-invariant
to the flip, so $V$---and therefore $\varphi$---barely moves, and the AUROC ceiling is
$\approx 0.5$ \emph{by construction}, not through any deficiency of the detector. The small,
consistent residual below chance (a few points) reflects imperfect symmetry, i.e.\ a faint
directional shift rather than noise. The honest conclusion is that polarity\_flip is largely
undetectable from the membrane and its signal lives on the \emph{input} side (the
positive/negative event-count balance the flip inverts); we scope it accordingly rather than
chase it.

\subsection{Why GAP causes the inversions, and what the theory prescribes}
\label{sec:theory-taxonomy}

The three hard corruptions---spatial\_dropout, event\_flood, polarity\_flip---share a single
mechanism in the analysis above: their discriminative signature lives in structure that
GAP-then-moments discards (spatial dispersion for the first two, input-side sign balance for
the third), while the surviving marginal moments are preserved or even pushed toward the
clean mean. This yields a clean, falsifiable taxonomy that organizes the entire results
section:
%
\begin{itemize}
  \item \textbf{Magnitude corruptions} (hot\_pixel, event\_rate\_shift, event\_flood-as-mean,
        temporal\_jitter via variance) move the membrane moments \emph{away} from clean and
        are caught by distance-based detectors on $\varphi$.
  \item \textbf{Structure corruptions} (spatial\_dropout, and the spatial component of
        event\_flood) leave the marginal moments intact and require either the
        spatial-dispersion features GAP discards or a scoring rule that flags ``too
        typical,'' not only ``too far.''
  \item \textbf{Invariant corruptions} (polarity\_flip) are near-orthogonal to the membrane
        and are an input-side, not a membrane-side, detection problem.
\end{itemize}
%
The theory therefore does three things at once: it \emph{explains} the easy wins as
inevitable, it \emph{predicts} the inversions from first principles, and it \emph{prescribes}
the architecture that addresses them---the spatial branch, two-sided scoring, and sequence
aggregation that the MDD detector (Section~\ref{sec:method}) combines. Crucially, it does so
without any free parameters fit to the corruption labels: every prediction follows from the
PLIF state equation and the definition of the feature.

\section{The Neuromorphic Hardware Advantage}
\label{sec:hardware}

The empirical case for membrane-statistics OOD detection is incomplete without its central
positioning claim, which is not about accuracy at all. Existing OOD detectors are designed
for artificial neural networks executed on von Neumann hardware, where logits, gradients, and
repeated forward passes are cheap. Spiking networks are deployed for the opposite reason---to
run within the severe energy and compute budget of neuromorphic silicon---and in that regime
the standard OOD toolkit is not merely expensive but \emph{architecturally infeasible}.
Membrane statistics are the rare signal that is free precisely where every alternative is
impossible. This reframes the contribution from ``a competitive OOD score'' to ``the only OOD
method that runs on the hardware the model targets.''

\subsection{The membrane potential is a native, already-maintained register}
\label{sec:hardware-register}

On the major neuromorphic platforms---Intel Loihi \cite{davies2018loihi}, the BrainScaleS
analog system \cite{schemmel2010brainscales}, and SpiNNaker \cite{furber2014spinnaker}---each
neuron's membrane potential $V$ is \emph{physical state the substrate already maintains}. It
is not an intermediate that must be recomputed for our purposes; it is the quantity the neuron
updates on every time step in order to decide whether to spike, held in a local register
(digital cores) or a capacitor voltage (analog). Producing our descriptor $\varphi$ from it
requires only a read of that state and a reduction to three per-channel moments---an
$O(\text{channels})$ operation with no additional synaptic integration. In the cost accounting
that matters on-chip---multiply--accumulate operations and memory traffic for weights---%
$\varphi$ adds \emph{zero MACs and zero weight fetches}: it is a byproduct of the forward pass
that has to happen anyway to emit spikes. This is the precise sense in which the method is
``zero additional compute,'' and it is a property of \emph{where} the signal comes from, not
an optimization we apply afterward.

\subsection{Why competing OOD methods cannot run on the target hardware}
\label{sec:hardware-table}

The dominant OOD families each presuppose a capability that the neuromorphic execution model
does not provide. We enumerate them not to claim our accuracy is higher, but to establish that
the comparison cannot even take place on-chip.

\begin{table}[t]
  \centering
  \small
  \begin{tabularx}{\linewidth}{@{}l X X@{}}
    \toprule
    Method & Required extra computation & Why infeasible on neuromorphic HW \\
    \midrule
    \textbf{Vmem-$\varphi$ (ours)} & native $V$ register read + moments & --- (it is free) \\
    MSP \cite{hendrycks2017msp} / Energy \cite{liu2020energy}
      & softmax / log-sum-exp over logits
      & cores emit spikes, not float logits; no on-chip softmax \\
    ODIN \cite{liang2018odin}
      & temperature softmax + input perturbation
      & perturbation needs an input-gradient pass the chip cannot take \\
    ReAct \cite{sun2021react}
      & activation clipping + second forward pass
      & re-forwarding violates the SNN energy budget \\
    ViM \cite{wang2022vim}
      & SVD / dense linear algebra on features
      & no dense float linear-algebra unit; data movement alone too costly \\
    GradNorm \cite{huang2021gradnorm}
      & gradient backpropagation at inference
      & neuromorphic inference has no backward pass \\
    kNN / Mahalanobis on \emph{ANN} features \cite{lee2018mahalanobis,sun2022knn}
      & an auxiliary ANN forward pass
      & requires a second, non-spiking network---defeating the SNN premise \\
    \bottomrule
  \end{tabularx}
  \caption{Standard OOD detectors each require a capability absent from the spiking,
  feed-forward, integer-event execution model; Vmem-$\varphi$ requires none because the signal
  is the neuron state itself.}
  \label{tab:hardware}
\end{table}

The pattern is uniform: every alternative needs either a floating-point readout layer
(softmax, SVD), a second pass (ReAct, ANN-feature methods), or a gradient (ODIN,
GradNorm)---none of which exist in the spiking, feed-forward, integer-event execution model.
Vmem-$\varphi$ needs none of them because the signal is the neuron state itself.

\subsection{The reframed contribution}
\label{sec:hardware-claim}

It follows that the benchmark of Section~\ref{sec:experiments} should be read not as
``$\varphi$ narrowly beats ANN OOD baselines on a GPU'' but as ``$\varphi$ provides usable OOD
detection \emph{inside the compute envelope of the deployment hardware}, where the listed
alternatives cannot be evaluated at all.'' This is a categorical claim about feasibility, and
it is what makes the work a distinct contribution rather than another row in an OOD
leaderboard: for spiking networks on neuromorphic silicon, membrane-statistics OOD detection
is not the best option among many---under a strict on-chip compute budget it is, to our
knowledge, the only option that is free.

\subsection{Scope and honesty caveat}
\label{sec:hardware-caveat}

We measure $\varphi$ in simulation (SpikingJelly on a GPU), not on neuromorphic silicon. Our
cost claim is therefore \emph{architectural}---it follows from $V$ being a register the
substrate already maintains and from $\varphi$ adding no synaptic operations---rather than a
wall-clock or joule measurement on-chip. We state this explicitly: the contribution is that
the signal is free \emph{by construction} on the target hardware, and an on-silicon energy
measurement is identified as future work. We likewise restrict the infeasibility argument to
the standard execution models of current neuromorphic platforms; co-processor or hybrid
configurations could in principle host some of the listed methods at a cost, which only
sharpens the comparison rather than changing its direction.

% --- Standalone bibliography (skipped when embedded). Manual thebibliography so
%     the standalone build needs no .bib file and no bibtex pass. Replace with a
%     real \bibliography{...} when merging into the full manuscript.
\ifdefined\theoryhardwareEMBED\else
\begin{thebibliography}{99}

\bibitem{fang2021plif}
W.~Fang, Z.~Yu, Y.~Chen, T.~Masquelier, T.~Huang, and Y.~Tian.
\newblock Incorporating learnable membrane time constant to enhance learning of
  spiking neural networks.
\newblock In \emph{ICCV}, 2021.

\bibitem{davies2018loihi}
M.~Davies et~al.
\newblock Loihi: A neuromorphic manycore processor with on-chip learning.
\newblock \emph{IEEE Micro}, 38(1):82--99, 2018.

\bibitem{schemmel2010brainscales}
J.~Schemmel, D.~Br\"uderle, A.~Gr\"ubl, M.~Hock, K.~Meier, and S.~Millner.
\newblock A wafer-scale neuromorphic hardware system for large-scale neural
  modeling.
\newblock In \emph{ISCAS}, pages 1947--1950, 2010.

\bibitem{furber2014spinnaker}
S.~B. Furber, F.~Galluppi, S.~Temple, and L.~A. Plana.
\newblock The SpiNNaker project.
\newblock \emph{Proceedings of the IEEE}, 102(5):652--665, 2014.

\bibitem{hendrycks2017msp}
D.~Hendrycks and K.~Gimpel.
\newblock A baseline for detecting misclassified and out-of-distribution
  examples in neural networks.
\newblock In \emph{ICLR}, 2017.

\bibitem{liu2020energy}
W.~Liu, X.~Wang, J.~D. Owens, and Y.~Li.
\newblock Energy-based out-of-distribution detection.
\newblock In \emph{NeurIPS}, 2020.

\bibitem{liang2018odin}
S.~Liang, Y.~Li, and R.~Srikant.
\newblock Enhancing the reliability of out-of-distribution image detection in
  neural networks.
\newblock In \emph{ICLR}, 2018.

\bibitem{sun2021react}
Y.~Sun, C.~Guo, and Y.~Li.
\newblock ReAct: Out-of-distribution detection with rectified activations.
\newblock In \emph{NeurIPS}, 2021.

\bibitem{wang2022vim}
H.~Wang, Z.~Li, L.~Feng, and W.~Zhang.
\newblock ViM: Out-of-distribution with virtual-logit matching.
\newblock In \emph{CVPR}, 2022.

\bibitem{huang2021gradnorm}
R.~Huang, A.~Geng, and Y.~Li.
\newblock On the importance of gradients for detecting distributional shifts in
  the wild.
\newblock In \emph{NeurIPS}, 2021.

\bibitem{lee2018mahalanobis}
K.~Lee, K.~Lee, H.~Lee, and J.~Shin.
\newblock A simple unified framework for detecting out-of-distribution samples
  and adversarial attacks.
\newblock In \emph{NeurIPS}, 2018.

\bibitem{sun2022knn}
Y.~Sun, Y.~Ming, X.~Zhu, and Y.~Li.
\newblock Out-of-distribution detection with deep nearest neighbors.
\newblock In \emph{ICML}, 2022.

\end{thebibliography}
\end{document}
\fi

%
% Citation keys are placeholders; the embedded thebibliography below gives them
% resolvable (if approximate) entries so the standalone build is clean.

\ifdefined\theoryhardwareEMBED\else
\documentclass[11pt,a4paper]{article}
\usepackage[utf8]{inputenc}
\usepackage[margin=1in]{geometry}
\usepackage{amsmath}
\usepackage{amssymb}
\usepackage{booktabs}
\usepackage{tabularx}
\usepackage{microtype}
\begin{document}
\fi

\section{The Event-Camera Corruption Suite}
\label{sec:corruptions}

A benchmark for out-of-distribution (OOD) detection is only as meaningful as the
distribution shifts it tests against. We therefore define a controlled, reproducible
battery of six corruptions that model the dominant failure modes of event-camera
front-ends, each applied at five monotonically increasing severities. Together with the
clean baseline this yields $6\times5+1=31$ runs per detector. The suite is designed so that
each corruption perturbs a \emph{structurally distinct} axis of the input---DC offset,
localized burst, temporal order, global rate, polarity, or spatial support---which is what
makes the closed-form analysis of Section~\ref{sec:theory} possible: a detector's behaviour
on each corruption can be predicted from which axis it moves.

\subsection{Input representation}
\label{sec:corruptions-input}

Every corruption operates directly on the model's input tensor rather than on the raw event
stream, so that it composes with any downstream architecture and is label-free (no detection
annotations are consulted). A sequence is a stack of $N$ event-histogram frames of shape
$(20, H, W)$, with $H\times W = 240\times304$ for Prophesee Gen1. The $20$ channels encode
two polarities $\times$ ten time bins of width $50\,\mu\text{s}$: channels $0$--$9$ are the
ON-polarity bins (oldest to newest) and channels $10$--$19$ the OFF-polarity bins. Counts
are stored as \texttt{uint8}; corruptions that can increase counts promote to a wider integer
or float type internally and clip back to $[0,255]$, so no silent wrap-around occurs.

\subsection{Design principles}
\label{sec:corruptions-design}

Three properties hold across the suite. \textbf{(i) Reproducibility:} every random decision
(pixel coordinates, burst centres, flip masks, shift amounts, rate direction) is drawn from a
single seeded NumPy generator, so a (corruption, severity, seed) triple is deterministic.
\textbf{(ii) Hardware agnosticism:} each operation is expressed through an array module that
dispatches to NumPy or CuPy, so the identical corruption runs on CPU or GPU.
\textbf{(iii) Severity monotonicity:} severity $1$ is a mild perturbation and severity $5$ a
severe one, with parameters chosen to increase monotonically (Section~\ref{sec:corr-algos}),
enabling the Spearman-$\rho$ monotonicity analysis of the results section. Table~\ref{tab:corr-taxonomy}
summarizes the failure mode and perturbed axis for each corruption.

\begin{table}[t]
  \centering
  \small
  \begin{tabularx}{\linewidth}{@{}l l X@{}}
    \toprule
    Corruption & Physical failure mode & Axis of the input perturbed \\
    \midrule
    \texttt{hot\_pixel}        & defective / stuck pixels        & constant DC offset at fixed pixel locations \\
    \texttt{event\_flood}      & EMI / sensor saturation         & localized count bursts in space and time \\
    \texttt{temporal\_jitter}  & clock noise / timestamp error   & time-bin ordering (per polarity) \\
    \texttt{event\_rate\_shift}& event-density domain gap         & global multiplicative rate \\
    \texttt{polarity\_flip}    & ON/OFF misclassification         & polarity sign on a fraction of frames \\
    \texttt{spatial\_dropout}  & dead-pixel clusters / occlusion & spatial support (rectangular regions zeroed) \\
    \bottomrule
  \end{tabularx}
  \caption{The six event-camera corruptions, the sensor failure each models, and the
  structural axis of the input histogram it perturbs. Each axis maps to a distinct,
  predictable signature in the membrane moments (Section~\ref{sec:theory}).}
  \label{tab:corr-taxonomy}
\end{table}

\section{Corruption Algorithms}
\label{sec:corr-algos}

We now specify each corruption and its severity schedule. All operate on a histogram
$X \in \{0,\dots,255\}^{N\times 20\times H\times W}$ and return a corrupted copy of the same
shape and dtype.

\paragraph{Hot pixel.}
A fixed set of $n$ pixel locations is sampled uniformly at random and a constant count $\delta$
is added to \emph{every} time bin of \emph{every} frame at those locations, modelling pixels
that fire continuously regardless of scene content. The accumulation is performed in
\texttt{int16} and clipped to $[0,255]$. Severity raises both the number of hot pixels and the
injected count: $(n,\delta) \in \{(10,20),(30,40),(80,80),(150,140),(300,200)\}$ for
severities $1$--$5$.

\paragraph{Event flood.}
$b$ burst centres are placed at random frames and random spatial patches; at each, a count
$\delta$ is added inside a patch covering a fraction $f$ of the array, over the frames within
$\pm 2$ of the centre. This models electromagnetic interference or saturation affecting a
localized region for a short time. Severity raises the number of bursts, the patch fraction,
and the injected count: $(b,f,\delta) \in \{(2,0.05,60),(4,0.08,100),(8,0.12,150),(15,0.18,200),(25,0.25,255)\}$.

\paragraph{Temporal jitter.}
The ON-polarity bin block (channels $0$--$9$) and the OFF block (channels $10$--$19$) are each
rolled along the time-bin axis by an independent random integer shift drawn from
$[-s_{\max}, s_{\max}]$, with vacated bins zero-filled. This models clock noise / timestamp
inaccuracy that misassigns events to neighbouring time bins; crucially it permutes time order
while leaving the per-channel marginal counts almost unchanged. The maximum shift grows with
severity: $s_{\max} \in \{1,2,3,5,8\}$ bins (one bin $=50\,\mu\text{s}$).

\paragraph{Event-rate shift.}
All counts are multiplied by a single factor and clipped, modelling a domain gap in event
density (low- vs.\ high-texture or slow vs.\ fast motion). A seeded coin flip selects an
under-rate or over-rate factor, so the corruption is two-sided. The factor pair
$(\text{under},\text{over})$ widens with severity:
$\{(0.80,1.20),(0.65,1.40),(0.50,1.65),(0.35,2.00),(0.20,2.50)\}$.

\paragraph{Polarity flip.}
For a random fraction $p$ of frames, the ON and OFF channel blocks are swapped, modelling
sensor misclassification of brightness-change direction. The fraction of affected frames grows
with severity: $p \in \{0.05,0.10,0.20,0.35,0.50\}$.

\paragraph{Spatial dropout.}
$k$ rectangular regions, each covering fractions $(f_h, f_w)$ of the height and width, are
zeroed across all channels and all frames, modelling dead-pixel clusters or persistent
structured occlusion. Severity raises the number and size of regions:
$(k, f_h, f_w) \in \{(1,0.05,0.05),(1,0.10,0.10),(2,0.10,0.10),(2,0.15,0.15),(3,0.20,0.20)\}$.

\begin{table}[t]
  \centering
  \small
  \begin{tabularx}{\linewidth}{@{}l X X X X X@{}}
    \toprule
    Corruption & Sev.\ 1 & Sev.\ 2 & Sev.\ 3 & Sev.\ 4 & Sev.\ 5 \\
    \midrule
    \texttt{hot\_pixel} $(n,\delta)$        & $10,20$  & $30,40$  & $80,80$   & $150,140$ & $300,200$ \\
    \texttt{event\_flood} $(b,f,\delta)$    & $2,.05,60$ & $4,.08,100$ & $8,.12,150$ & $15,.18,200$ & $25,.25,255$ \\
    \texttt{temporal\_jitter} $s_{\max}$    & $1$ & $2$ & $3$ & $5$ & $8$ \\
    \texttt{event\_rate\_shift} (u, o)      & $.80,1.20$ & $.65,1.40$ & $.50,1.65$ & $.35,2.00$ & $.20,2.50$ \\
    \texttt{polarity\_flip} $p$             & $.05$ & $.10$ & $.20$ & $.35$ & $.50$ \\
    \texttt{spatial\_dropout} $(k,f_h,f_w)$ & $1,.05,.05$ & $1,.10,.10$ & $2,.10,.10$ & $2,.15,.15$ & $3,.20,.20$ \\
    \bottomrule
  \end{tabularx}
  \caption{Severity schedules for the six corruptions. Parameters increase monotonically with
  severity. Symbols: $n$ hot pixels, $\delta$ injected count, $b$ burst centres, $f$ patch
  fraction, $s_{\max}$ max bin shift, $(u,o)$ under/over rate factors, $p$ flipped-frame
  fraction, $k$ dropout regions, $(f_h,f_w)$ region height/width fractions.}
  \label{tab:corr-severity}
\end{table}

\section{A Leaky-Integrator Theory of Membrane-Statistics OOD Detection}
\label{sec:theory}

A benchmark that reports which corruptions a detector catches, but not \emph{why},
invites the obvious reviewer objection: the method may be exploiting dataset-specific
artifacts rather than a property of the underlying neuron model. This section removes that
objection. We show that the detectability of each corruption is not an empirical accident
but a closed-form consequence of the Parametric Leaky Integrate-and-Fire (PLIF) state
equation. The membrane potential is a \emph{leaky exponential accumulator} of the input
current; each corruption perturbs that current in a structurally distinct way, and the
perturbation propagates to a predictable change in the first three moments of the
sub-threshold trace. The resulting predictions match the measured per-corruption AUROC
ordering, turning the benchmark into a principled framework.

\subsection{The PLIF state equation and its closed-form solution}
\label{sec:theory-state}

The PLIF neuron, as instantiated by SpikingJelly's \texttt{MultiStepParametricLIFNode}
\cite{fang2021plif}, maintains a sub-threshold membrane potential $V$ that integrates an
input current with a learnable leak. In discrete multi-step form, the update for a single
channel is
%
\begin{equation}
  V[t] \;=\; \Big(1 - \tfrac{1}{\tau}\Big)\, V[t-1] \;+\; W\, X[t],
  \label{eq:plif}
\end{equation}
%
where $\tau > 1$ is the membrane time constant (learned per layer), $W$ is the synaptic
weight applied to the input event-histogram current $X[t]$ at time step $t$, and a spike is
emitted with a reset whenever $V$ crosses the firing threshold $\theta$. Writing the leak
factor as $\lambda \equiv 1 - 1/\tau \in (0,1)$ and unrolling \eqref{eq:plif} from a
quiescent initial state $V[-1]=0$ over the $T$ SNN time steps, the sub-threshold trace is an
exponentially weighted history of the input:
%
\begin{equation}
  V[t] \;=\; \sum_{k=0}^{t} \lambda^{\,k}\, W\, X[t-k].
  \label{eq:unroll}
\end{equation}
%
Equation~\eqref{eq:unroll} is the object the entire method rests on: $V$ is a \emph{causal,
leaky linear filter} of the event stream, with an effective memory horizon of $\sim\tau$
steps. Two immediate consequences frame everything that follows. First, $V$ is \emph{not} a
memoryless function of the current frame---it carries the temporal structure of recent
input, which is why temporally-structured corruptions leave a signature even when the
marginal input distribution is unchanged. Second, because the map from $X$ to $V$ is linear,
the \emph{moments} of $V$ inherit closed-form dependence on the moments of $X$, which we now
make explicit.

\subsection{The feature: per-channel membrane moments}
\label{sec:theory-feature}

For each channel $c$ in each of the four PLIF layers, we apply Global Average Pooling (GAP)
over the spatial dimensions of $V$ and compute the first three central moments of the pooled
trace across the $T$ steps:
%
\begin{equation}
  \varphi_c \;=\; \big[\, \mu_c,\; \sigma_c^2,\; \kappa_c \,\big],
  \qquad
  \mu_c = \mathbb{E}[V_c],\quad
  \sigma_c^2 = \mathrm{Var}[V_c],\quad
  \kappa_c = \frac{\mathbb{E}\big[(V_c-\mu_c)^4\big]}{\sigma_c^4}.
  \label{eq:phi}
\end{equation}
%
Stacking over $704$ channels across the four layers yields the $2112$-dimensional descriptor
$\varphi$. The mean captures sustained drive, the variance captures fluctuation energy, and
the kurtosis captures burstiness (heavy-tailed, intermittent excitation). We show next that
these three coordinates are individually informative about \emph{different} corruption
families.

\subsection{Steady-state moments under a stationary input}
\label{sec:theory-moments}

Suppose the per-channel input current is wide-sense stationary with mean $m=\mathbb{E}[X]$,
variance $s^2=\mathrm{Var}[X]$, and autocovariance $\gamma(j)=\mathrm{Cov}(X[t],X[t-j])$.
Taking expectations of the geometric sum \eqref{eq:unroll} in the long-horizon limit
($\sum_k \lambda^k \to 1/(1-\lambda)$) gives the steady-state membrane \emph{mean}
%
\begin{equation}
  \mathbb{E}[V] \;=\; \frac{W m}{1-\lambda} \;=\; \tau\, W m,
  \label{eq:mean}
\end{equation}
%
so the mean is the input rate \emph{amplified by the time constant} $\tau$. The steady-state
\emph{variance}, assuming weak temporal correlation ($\gamma(j)\!\approx\!0$ for $j>0$), is
%
\begin{equation}
  \mathrm{Var}[V] \;=\; (W s)^2 \sum_{k=0}^{\infty} \lambda^{2k}
  \;=\; \frac{(W s)^2}{1-\lambda^2}
  \;=\; \frac{\tau^2}{2\tau-1}\,(W s)^2,
  \label{eq:var}
\end{equation}
%
i.e.\ the membrane variance tracks the \emph{input variance}, scaled by a leak-dependent
constant. Retaining the autocovariance term generalizes \eqref{eq:var} to
%
\begin{equation}
  \mathrm{Var}[V] \;=\; \frac{(W s)^2}{1-\lambda^2}
  \;+\; 2 W^2 \sum_{j=1}^{\infty} \gamma(j)\,\frac{\lambda^{\,j}}{1-\lambda^2},
  \label{eq:var-auto}
\end{equation}
%
which is the key to temporal corruptions: the second term depends on the \emph{time
ordering} of the input, not its marginal distribution. A corruption that preserves $m$ and
$s^2$ but scrambles $\gamma(j)$ moves $\mathrm{Var}[V]$ (and, through higher-order analogues,
the kurtosis and the autocorrelation of $V(t)$) while leaving the naive first-moment picture
untouched. Equations~\eqref{eq:mean}--\eqref{eq:var-auto} are the three levers every
corruption pulls; the rest of the section reads off, corruption by corruption, which lever
moves and therefore which feature detects it.

\subsection{Per-corruption analysis}
\label{sec:theory-percorruption}

\paragraph{Hot-pixel --- DC saturation (predicted AUROC $\approx 1$).}
A stuck/hot pixel injects a persistent constant $\delta$ into $X[t]$ at fixed spatial
locations for the affected channels. Adding a constant $\delta$ to the input in
\eqref{eq:mean} shifts the membrane mean by $\tau W \delta$, an unbounded accumulation that
drives $V$ monotonically toward---and clamps it at---the firing threshold $\theta$. The
affected channels' $\mu_c$ therefore saturate at a value disjoint from the clean range,
producing two non-overlapping clusters in $\varphi$-space. Detection is not merely easy but
\emph{mathematically inevitable}, and it is monotone in severity because larger $\delta$
saturates more channels. This is exactly the measured behaviour (perfect AUROC and Spearman
monotonicity).

\paragraph{Temporal jitter --- autocorrelation destruction (temporal/deep features only).}
Temporal jitter permutes the time order of $X[t]$ within a window. A permutation is
measure-preserving: it leaves the marginal mean $m$ and variance $s^2$---and hence
\eqref{eq:mean} and the \emph{first} term of \eqref{eq:var-auto}---unchanged. What it
destroys is the autocovariance $\gamma(j)$, i.e.\ the second term of \eqref{eq:var-auto} and,
more sharply, the lag-1 autocorrelation of the membrane trace $V(t)$ itself. The leaky filter
\eqref{eq:unroll} is precisely an autocorrelation-generating operation, so jitter is visible
in the \emph{temporal} structure of $V$ even though it is invisible to the static marginal
mean. This predicts, and the data confirm, that jitter is caught by temporal statistics
(lag-1 autocorrelation) and by deeper layers whose larger effective $\tau$ integrate
longer-range temporal structure---the layer-4 slice is the strongest single view for jitter.

\paragraph{Event-rate shift --- global affine drive (single activity scalar).}
An event-rate shift multiplies the global input rate $m \to \alpha m$. By \eqref{eq:mean}
this scales the membrane mean of \emph{every} channel by the same factor $\alpha$, an
approximately affine displacement of $\varphi$ along a single global-activity direction. The
corruption is thus detectable by a one-dimensional projection---a global activity
scalar---which is why a single summary statistic recovers most of the signal, and why the
per-frame AUROC shows a phase transition once $\alpha$ leaves the clean operating band.

\paragraph{Spatial dropout --- variance loss hidden by GAP (anti-detectable under one-sided
scoring).}
Spatial dropout zeros a contiguous spatial \emph{region} of $X[t]$. Within an affected
channel this lowers the spatial input variance and, through \eqref{eq:var}, the membrane
variance $\sigma_c^2$, while leaving the channel mean $\mu_c$ comparatively intact (the
surviving region still drives the neuron). Variance is therefore the single most informative
moment for this corruption---but GAP pools away the \emph{spatial} dimension before the
moments are taken, so the dominant signature (which spatial regions went silent) is
discarded, and what remains makes the membrane look \emph{quieter and more regular}, i.e.\
closer to the clean mean. Under one-sided distance-from-clean scoring this inverts the
detector (AUROC $< 0.5$): the corruption is anti-detectable not because the information is
absent from $V$ but because GAP plus a one-sided rule throws it away. This is the formal
motivation for the spatial-dispersion features (\texttt{phi\_spatial}) and for two-sided
scoring (Section~\ref{sec:method}).

\paragraph{Event flood --- uniform mean shift indistinguishable from a busy scene
(near-chance per frame).}
Event flood adds approximately uniform noise across all pixels, raising $m$ uniformly. By
\eqref{eq:mean} this lifts the mean of \emph{all} channels by nearly the same amount---the
identical signature a genuinely busy, high-activity clean scene produces. Because the clean
distribution already contains such high-activity scenes, the flooded frames are not separable
from the clean tail on the marginal moments alone, giving near-chance per-frame AUROC. The
information that \emph{does} distinguish flood---that the elevated activity is spatially
unstructured (it fills space rather than tracking objects)---again lives in the spatial
dispersion GAP removes, and in the temporal \emph{consistency} of the bias across a sequence.
This predicts the two recoveries we observe: spatial features, and sequence-level
aggregation, which lifts flood to near-perfect AUROC once the per-frame bias is integrated
over a whole recording.

\paragraph{Polarity flip --- invariance by construction ($\approx 0.5$ ceiling, a
membrane-side dead end).}
Polarity flip inverts the ON/OFF sign of the events. The upstream detector was trained to be
approximately polarity-symmetric; to the extent that symmetry holds, $W X$ is near-invariant
to the flip, so $V$---and therefore $\varphi$---barely moves, and the AUROC ceiling is
$\approx 0.5$ \emph{by construction}, not through any deficiency of the detector. The small,
consistent residual below chance (a few points) reflects imperfect symmetry, i.e.\ a faint
directional shift rather than noise. The honest conclusion is that polarity\_flip is largely
undetectable from the membrane and its signal lives on the \emph{input} side (the
positive/negative event-count balance the flip inverts); we scope it accordingly rather than
chase it.

\subsection{Why GAP causes the inversions, and what the theory prescribes}
\label{sec:theory-taxonomy}

The three hard corruptions---spatial\_dropout, event\_flood, polarity\_flip---share a single
mechanism in the analysis above: their discriminative signature lives in structure that
GAP-then-moments discards (spatial dispersion for the first two, input-side sign balance for
the third), while the surviving marginal moments are preserved or even pushed toward the
clean mean. This yields a clean, falsifiable taxonomy that organizes the entire results
section:
%
\begin{itemize}
  \item \textbf{Magnitude corruptions} (hot\_pixel, event\_rate\_shift, event\_flood-as-mean,
        temporal\_jitter via variance) move the membrane moments \emph{away} from clean and
        are caught by distance-based detectors on $\varphi$.
  \item \textbf{Structure corruptions} (spatial\_dropout, and the spatial component of
        event\_flood) leave the marginal moments intact and require either the
        spatial-dispersion features GAP discards or a scoring rule that flags ``too
        typical,'' not only ``too far.''
  \item \textbf{Invariant corruptions} (polarity\_flip) are near-orthogonal to the membrane
        and are an input-side, not a membrane-side, detection problem.
\end{itemize}
%
The theory therefore does three things at once: it \emph{explains} the easy wins as
inevitable, it \emph{predicts} the inversions from first principles, and it \emph{prescribes}
the architecture that addresses them---the spatial branch, two-sided scoring, and sequence
aggregation that the MDD detector (Section~\ref{sec:method}) combines. Crucially, it does so
without any free parameters fit to the corruption labels: every prediction follows from the
PLIF state equation and the definition of the feature.

\section{The Neuromorphic Hardware Advantage}
\label{sec:hardware}

The empirical case for membrane-statistics OOD detection is incomplete without its central
positioning claim, which is not about accuracy at all. Existing OOD detectors are designed
for artificial neural networks executed on von Neumann hardware, where logits, gradients, and
repeated forward passes are cheap. Spiking networks are deployed for the opposite reason---to
run within the severe energy and compute budget of neuromorphic silicon---and in that regime
the standard OOD toolkit is not merely expensive but \emph{architecturally infeasible}.
Membrane statistics are the rare signal that is free precisely where every alternative is
impossible. This reframes the contribution from ``a competitive OOD score'' to ``the only OOD
method that runs on the hardware the model targets.''

\subsection{The membrane potential is a native, already-maintained register}
\label{sec:hardware-register}

On the major neuromorphic platforms---Intel Loihi \cite{davies2018loihi}, the BrainScaleS
analog system \cite{schemmel2010brainscales}, and SpiNNaker \cite{furber2014spinnaker}---each
neuron's membrane potential $V$ is \emph{physical state the substrate already maintains}. It
is not an intermediate that must be recomputed for our purposes; it is the quantity the neuron
updates on every time step in order to decide whether to spike, held in a local register
(digital cores) or a capacitor voltage (analog). Producing our descriptor $\varphi$ from it
requires only a read of that state and a reduction to three per-channel moments---an
$O(\text{channels})$ operation with no additional synaptic integration. In the cost accounting
that matters on-chip---multiply--accumulate operations and memory traffic for weights---%
$\varphi$ adds \emph{zero MACs and zero weight fetches}: it is a byproduct of the forward pass
that has to happen anyway to emit spikes. This is the precise sense in which the method is
``zero additional compute,'' and it is a property of \emph{where} the signal comes from, not
an optimization we apply afterward.

\subsection{Why competing OOD methods cannot run on the target hardware}
\label{sec:hardware-table}

The dominant OOD families each presuppose a capability that the neuromorphic execution model
does not provide. We enumerate them not to claim our accuracy is higher, but to establish that
the comparison cannot even take place on-chip.

\begin{table}[t]
  \centering
  \small
  \begin{tabularx}{\linewidth}{@{}l X X@{}}
    \toprule
    Method & Required extra computation & Why infeasible on neuromorphic HW \\
    \midrule
    \textbf{Vmem-$\varphi$ (ours)} & native $V$ register read + moments & --- (it is free) \\
    MSP \cite{hendrycks2017msp} / Energy \cite{liu2020energy}
      & softmax / log-sum-exp over logits
      & cores emit spikes, not float logits; no on-chip softmax \\
    ODIN \cite{liang2018odin}
      & temperature softmax + input perturbation
      & perturbation needs an input-gradient pass the chip cannot take \\
    ReAct \cite{sun2021react}
      & activation clipping + second forward pass
      & re-forwarding violates the SNN energy budget \\
    ViM \cite{wang2022vim}
      & SVD / dense linear algebra on features
      & no dense float linear-algebra unit; data movement alone too costly \\
    GradNorm \cite{huang2021gradnorm}
      & gradient backpropagation at inference
      & neuromorphic inference has no backward pass \\
    kNN / Mahalanobis on \emph{ANN} features \cite{lee2018mahalanobis,sun2022knn}
      & an auxiliary ANN forward pass
      & requires a second, non-spiking network---defeating the SNN premise \\
    \bottomrule
  \end{tabularx}
  \caption{Standard OOD detectors each require a capability absent from the spiking,
  feed-forward, integer-event execution model; Vmem-$\varphi$ requires none because the signal
  is the neuron state itself.}
  \label{tab:hardware}
\end{table}

The pattern is uniform: every alternative needs either a floating-point readout layer
(softmax, SVD), a second pass (ReAct, ANN-feature methods), or a gradient (ODIN,
GradNorm)---none of which exist in the spiking, feed-forward, integer-event execution model.
Vmem-$\varphi$ needs none of them because the signal is the neuron state itself.

\subsection{The reframed contribution}
\label{sec:hardware-claim}

It follows that the benchmark of Section~\ref{sec:experiments} should be read not as
``$\varphi$ narrowly beats ANN OOD baselines on a GPU'' but as ``$\varphi$ provides usable OOD
detection \emph{inside the compute envelope of the deployment hardware}, where the listed
alternatives cannot be evaluated at all.'' This is a categorical claim about feasibility, and
it is what makes the work a distinct contribution rather than another row in an OOD
leaderboard: for spiking networks on neuromorphic silicon, membrane-statistics OOD detection
is not the best option among many---under a strict on-chip compute budget it is, to our
knowledge, the only option that is free.

\subsection{Scope and honesty caveat}
\label{sec:hardware-caveat}

We measure $\varphi$ in simulation (SpikingJelly on a GPU), not on neuromorphic silicon. Our
cost claim is therefore \emph{architectural}---it follows from $V$ being a register the
substrate already maintains and from $\varphi$ adding no synaptic operations---rather than a
wall-clock or joule measurement on-chip. We state this explicitly: the contribution is that
the signal is free \emph{by construction} on the target hardware, and an on-silicon energy
measurement is identified as future work. We likewise restrict the infeasibility argument to
the standard execution models of current neuromorphic platforms; co-processor or hybrid
configurations could in principle host some of the listed methods at a cost, which only
sharpens the comparison rather than changing its direction.

% --- Standalone bibliography (skipped when embedded). Manual thebibliography so
%     the standalone build needs no .bib file and no bibtex pass. Replace with a
%     real \bibliography{...} when merging into the full manuscript.
\ifdefined\theoryhardwareEMBED\else
\begin{thebibliography}{99}

\bibitem{fang2021plif}
W.~Fang, Z.~Yu, Y.~Chen, T.~Masquelier, T.~Huang, and Y.~Tian.
\newblock Incorporating learnable membrane time constant to enhance learning of
  spiking neural networks.
\newblock In \emph{ICCV}, 2021.

\bibitem{davies2018loihi}
M.~Davies et~al.
\newblock Loihi: A neuromorphic manycore processor with on-chip learning.
\newblock \emph{IEEE Micro}, 38(1):82--99, 2018.

\bibitem{schemmel2010brainscales}
J.~Schemmel, D.~Br\"uderle, A.~Gr\"ubl, M.~Hock, K.~Meier, and S.~Millner.
\newblock A wafer-scale neuromorphic hardware system for large-scale neural
  modeling.
\newblock In \emph{ISCAS}, pages 1947--1950, 2010.

\bibitem{furber2014spinnaker}
S.~B. Furber, F.~Galluppi, S.~Temple, and L.~A. Plana.
\newblock The SpiNNaker project.
\newblock \emph{Proceedings of the IEEE}, 102(5):652--665, 2014.

\bibitem{hendrycks2017msp}
D.~Hendrycks and K.~Gimpel.
\newblock A baseline for detecting misclassified and out-of-distribution
  examples in neural networks.
\newblock In \emph{ICLR}, 2017.

\bibitem{liu2020energy}
W.~Liu, X.~Wang, J.~D. Owens, and Y.~Li.
\newblock Energy-based out-of-distribution detection.
\newblock In \emph{NeurIPS}, 2020.

\bibitem{liang2018odin}
S.~Liang, Y.~Li, and R.~Srikant.
\newblock Enhancing the reliability of out-of-distribution image detection in
  neural networks.
\newblock In \emph{ICLR}, 2018.

\bibitem{sun2021react}
Y.~Sun, C.~Guo, and Y.~Li.
\newblock ReAct: Out-of-distribution detection with rectified activations.
\newblock In \emph{NeurIPS}, 2021.

\bibitem{wang2022vim}
H.~Wang, Z.~Li, L.~Feng, and W.~Zhang.
\newblock ViM: Out-of-distribution with virtual-logit matching.
\newblock In \emph{CVPR}, 2022.

\bibitem{huang2021gradnorm}
R.~Huang, A.~Geng, and Y.~Li.
\newblock On the importance of gradients for detecting distributional shifts in
  the wild.
\newblock In \emph{NeurIPS}, 2021.

\bibitem{lee2018mahalanobis}
K.~Lee, K.~Lee, H.~Lee, and J.~Shin.
\newblock A simple unified framework for detecting out-of-distribution samples
  and adversarial attacks.
\newblock In \emph{NeurIPS}, 2018.

\bibitem{sun2022knn}
Y.~Sun, Y.~Ming, X.~Zhu, and Y.~Li.
\newblock Out-of-distribution detection with deep nearest neighbors.
\newblock In \emph{ICML}, 2022.

\end{thebibliography}
\end{document}
\fi

%
% Citation keys are placeholders; the embedded thebibliography below gives them
% resolvable (if approximate) entries so the standalone build is clean.

\ifdefined\theoryhardwareEMBED\else
\documentclass[11pt,a4paper]{article}
\usepackage[utf8]{inputenc}
\usepackage[margin=1in]{geometry}
\usepackage{amsmath}
\usepackage{amssymb}
\usepackage{booktabs}
\usepackage{tabularx}
\usepackage{microtype}
\begin{document}
\fi

\section{The Event-Camera Corruption Suite}
\label{sec:corruptions}

A benchmark for out-of-distribution (OOD) detection is only as meaningful as the
distribution shifts it tests against. We therefore define a controlled, reproducible
battery of six corruptions that model the dominant failure modes of event-camera
front-ends, each applied at five monotonically increasing severities. Together with the
clean baseline this yields $6\times5+1=31$ runs per detector. The suite is designed so that
each corruption perturbs a \emph{structurally distinct} axis of the input---DC offset,
localized burst, temporal order, global rate, polarity, or spatial support---which is what
makes the closed-form analysis of Section~\ref{sec:theory} possible: a detector's behaviour
on each corruption can be predicted from which axis it moves.

\subsection{Input representation}
\label{sec:corruptions-input}

Every corruption operates directly on the model's input tensor rather than on the raw event
stream, so that it composes with any downstream architecture and is label-free (no detection
annotations are consulted). A sequence is a stack of $N$ event-histogram frames of shape
$(20, H, W)$, with $H\times W = 240\times304$ for Prophesee Gen1. The $20$ channels encode
two polarities $\times$ ten time bins of width $50\,\mu\text{s}$: channels $0$--$9$ are the
ON-polarity bins (oldest to newest) and channels $10$--$19$ the OFF-polarity bins. Counts
are stored as \texttt{uint8}; corruptions that can increase counts promote to a wider integer
or float type internally and clip back to $[0,255]$, so no silent wrap-around occurs.

\subsection{Design principles}
\label{sec:corruptions-design}

Three properties hold across the suite. \textbf{(i) Reproducibility:} every random decision
(pixel coordinates, burst centres, flip masks, shift amounts, rate direction) is drawn from a
single seeded NumPy generator, so a (corruption, severity, seed) triple is deterministic.
\textbf{(ii) Hardware agnosticism:} each operation is expressed through an array module that
dispatches to NumPy or CuPy, so the identical corruption runs on CPU or GPU.
\textbf{(iii) Severity monotonicity:} severity $1$ is a mild perturbation and severity $5$ a
severe one, with parameters chosen to increase monotonically (Section~\ref{sec:corr-algos}),
enabling the Spearman-$\rho$ monotonicity analysis of the results section. Table~\ref{tab:corr-taxonomy}
summarizes the failure mode and perturbed axis for each corruption.

\begin{table}[t]
  \centering
  \small
  \begin{tabularx}{\linewidth}{@{}l l X@{}}
    \toprule
    Corruption & Physical failure mode & Axis of the input perturbed \\
    \midrule
    \texttt{hot\_pixel}        & defective / stuck pixels        & constant DC offset at fixed pixel locations \\
    \texttt{event\_flood}      & EMI / sensor saturation         & localized count bursts in space and time \\
    \texttt{temporal\_jitter}  & clock noise / timestamp error   & time-bin ordering (per polarity) \\
    \texttt{event\_rate\_shift}& event-density domain gap         & global multiplicative rate \\
    \texttt{polarity\_flip}    & ON/OFF misclassification         & polarity sign on a fraction of frames \\
    \texttt{spatial\_dropout}  & dead-pixel clusters / occlusion & spatial support (rectangular regions zeroed) \\
    \bottomrule
  \end{tabularx}
  \caption{The six event-camera corruptions, the sensor failure each models, and the
  structural axis of the input histogram it perturbs. Each axis maps to a distinct,
  predictable signature in the membrane moments (Section~\ref{sec:theory}).}
  \label{tab:corr-taxonomy}
\end{table}

\section{Corruption Algorithms}
\label{sec:corr-algos}

We now specify each corruption and its severity schedule. All operate on a histogram
$X \in \{0,\dots,255\}^{N\times 20\times H\times W}$ and return a corrupted copy of the same
shape and dtype.

\paragraph{Hot pixel.}
A fixed set of $n$ pixel locations is sampled uniformly at random and a constant count $\delta$
is added to \emph{every} time bin of \emph{every} frame at those locations, modelling pixels
that fire continuously regardless of scene content. The accumulation is performed in
\texttt{int16} and clipped to $[0,255]$. Severity raises both the number of hot pixels and the
injected count: $(n,\delta) \in \{(10,20),(30,40),(80,80),(150,140),(300,200)\}$ for
severities $1$--$5$.

\paragraph{Event flood.}
$b$ burst centres are placed at random frames and random spatial patches; at each, a count
$\delta$ is added inside a patch covering a fraction $f$ of the array, over the frames within
$\pm 2$ of the centre. This models electromagnetic interference or saturation affecting a
localized region for a short time. Severity raises the number of bursts, the patch fraction,
and the injected count: $(b,f,\delta) \in \{(2,0.05,60),(4,0.08,100),(8,0.12,150),(15,0.18,200),(25,0.25,255)\}$.

\paragraph{Temporal jitter.}
The ON-polarity bin block (channels $0$--$9$) and the OFF block (channels $10$--$19$) are each
rolled along the time-bin axis by an independent random integer shift drawn from
$[-s_{\max}, s_{\max}]$, with vacated bins zero-filled. This models clock noise / timestamp
inaccuracy that misassigns events to neighbouring time bins; crucially it permutes time order
while leaving the per-channel marginal counts almost unchanged. The maximum shift grows with
severity: $s_{\max} \in \{1,2,3,5,8\}$ bins (one bin $=50\,\mu\text{s}$).

\paragraph{Event-rate shift.}
All counts are multiplied by a single factor and clipped, modelling a domain gap in event
density (low- vs.\ high-texture or slow vs.\ fast motion). A seeded coin flip selects an
under-rate or over-rate factor, so the corruption is two-sided. The factor pair
$(\text{under},\text{over})$ widens with severity:
$\{(0.80,1.20),(0.65,1.40),(0.50,1.65),(0.35,2.00),(0.20,2.50)\}$.

\paragraph{Polarity flip.}
For a random fraction $p$ of frames, the ON and OFF channel blocks are swapped, modelling
sensor misclassification of brightness-change direction. The fraction of affected frames grows
with severity: $p \in \{0.05,0.10,0.20,0.35,0.50\}$.

\paragraph{Spatial dropout.}
$k$ rectangular regions, each covering fractions $(f_h, f_w)$ of the height and width, are
zeroed across all channels and all frames, modelling dead-pixel clusters or persistent
structured occlusion. Severity raises the number and size of regions:
$(k, f_h, f_w) \in \{(1,0.05,0.05),(1,0.10,0.10),(2,0.10,0.10),(2,0.15,0.15),(3,0.20,0.20)\}$.

\begin{table}[t]
  \centering
  \small
  \begin{tabularx}{\linewidth}{@{}l X X X X X@{}}
    \toprule
    Corruption & Sev.\ 1 & Sev.\ 2 & Sev.\ 3 & Sev.\ 4 & Sev.\ 5 \\
    \midrule
    \texttt{hot\_pixel} $(n,\delta)$        & $10,20$  & $30,40$  & $80,80$   & $150,140$ & $300,200$ \\
    \texttt{event\_flood} $(b,f,\delta)$    & $2,.05,60$ & $4,.08,100$ & $8,.12,150$ & $15,.18,200$ & $25,.25,255$ \\
    \texttt{temporal\_jitter} $s_{\max}$    & $1$ & $2$ & $3$ & $5$ & $8$ \\
    \texttt{event\_rate\_shift} (u, o)      & $.80,1.20$ & $.65,1.40$ & $.50,1.65$ & $.35,2.00$ & $.20,2.50$ \\
    \texttt{polarity\_flip} $p$             & $.05$ & $.10$ & $.20$ & $.35$ & $.50$ \\
    \texttt{spatial\_dropout} $(k,f_h,f_w)$ & $1,.05,.05$ & $1,.10,.10$ & $2,.10,.10$ & $2,.15,.15$ & $3,.20,.20$ \\
    \bottomrule
  \end{tabularx}
  \caption{Severity schedules for the six corruptions. Parameters increase monotonically with
  severity. Symbols: $n$ hot pixels, $\delta$ injected count, $b$ burst centres, $f$ patch
  fraction, $s_{\max}$ max bin shift, $(u,o)$ under/over rate factors, $p$ flipped-frame
  fraction, $k$ dropout regions, $(f_h,f_w)$ region height/width fractions.}
  \label{tab:corr-severity}
\end{table}

\section{A Leaky-Integrator Theory of Membrane-Statistics OOD Detection}
\label{sec:theory}

A benchmark that reports which corruptions a detector catches, but not \emph{why},
invites the obvious reviewer objection: the method may be exploiting dataset-specific
artifacts rather than a property of the underlying neuron model. This section removes that
objection. We show that the detectability of each corruption is not an empirical accident
but a closed-form consequence of the Parametric Leaky Integrate-and-Fire (PLIF) state
equation. The membrane potential is a \emph{leaky exponential accumulator} of the input
current; each corruption perturbs that current in a structurally distinct way, and the
perturbation propagates to a predictable change in the first three moments of the
sub-threshold trace. The resulting predictions match the measured per-corruption AUROC
ordering, turning the benchmark into a principled framework.

\subsection{The PLIF state equation and its closed-form solution}
\label{sec:theory-state}

The PLIF neuron, as instantiated by SpikingJelly's \texttt{MultiStepParametricLIFNode}
\cite{fang2021plif}, maintains a sub-threshold membrane potential $V$ that integrates an
input current with a learnable leak. In discrete multi-step form, the update for a single
channel is
%
\begin{equation}
  V[t] \;=\; \Big(1 - \tfrac{1}{\tau}\Big)\, V[t-1] \;+\; W\, X[t],
  \label{eq:plif}
\end{equation}
%
where $\tau > 1$ is the membrane time constant (learned per layer), $W$ is the synaptic
weight applied to the input event-histogram current $X[t]$ at time step $t$, and a spike is
emitted with a reset whenever $V$ crosses the firing threshold $\theta$. Writing the leak
factor as $\lambda \equiv 1 - 1/\tau \in (0,1)$ and unrolling \eqref{eq:plif} from a
quiescent initial state $V[-1]=0$ over the $T$ SNN time steps, the sub-threshold trace is an
exponentially weighted history of the input:
%
\begin{equation}
  V[t] \;=\; \sum_{k=0}^{t} \lambda^{\,k}\, W\, X[t-k].
  \label{eq:unroll}
\end{equation}
%
Equation~\eqref{eq:unroll} is the object the entire method rests on: $V$ is a \emph{causal,
leaky linear filter} of the event stream, with an effective memory horizon of $\sim\tau$
steps. Two immediate consequences frame everything that follows. First, $V$ is \emph{not} a
memoryless function of the current frame---it carries the temporal structure of recent
input, which is why temporally-structured corruptions leave a signature even when the
marginal input distribution is unchanged. Second, because the map from $X$ to $V$ is linear,
the \emph{moments} of $V$ inherit closed-form dependence on the moments of $X$, which we now
make explicit.

\subsection{The feature: per-channel membrane moments}
\label{sec:theory-feature}

For each channel $c$ in each of the four PLIF layers, we apply Global Average Pooling (GAP)
over the spatial dimensions of $V$ and compute the first three central moments of the pooled
trace across the $T$ steps:
%
\begin{equation}
  \varphi_c \;=\; \big[\, \mu_c,\; \sigma_c^2,\; \kappa_c \,\big],
  \qquad
  \mu_c = \mathbb{E}[V_c],\quad
  \sigma_c^2 = \mathrm{Var}[V_c],\quad
  \kappa_c = \frac{\mathbb{E}\big[(V_c-\mu_c)^4\big]}{\sigma_c^4}.
  \label{eq:phi}
\end{equation}
%
Stacking over $704$ channels across the four layers yields the $2112$-dimensional descriptor
$\varphi$. The mean captures sustained drive, the variance captures fluctuation energy, and
the kurtosis captures burstiness (heavy-tailed, intermittent excitation). We show next that
these three coordinates are individually informative about \emph{different} corruption
families.

\subsection{Steady-state moments under a stationary input}
\label{sec:theory-moments}

Suppose the per-channel input current is wide-sense stationary with mean $m=\mathbb{E}[X]$,
variance $s^2=\mathrm{Var}[X]$, and autocovariance $\gamma(j)=\mathrm{Cov}(X[t],X[t-j])$.
Taking expectations of the geometric sum \eqref{eq:unroll} in the long-horizon limit
($\sum_k \lambda^k \to 1/(1-\lambda)$) gives the steady-state membrane \emph{mean}
%
\begin{equation}
  \mathbb{E}[V] \;=\; \frac{W m}{1-\lambda} \;=\; \tau\, W m,
  \label{eq:mean}
\end{equation}
%
so the mean is the input rate \emph{amplified by the time constant} $\tau$. The steady-state
\emph{variance}, assuming weak temporal correlation ($\gamma(j)\!\approx\!0$ for $j>0$), is
%
\begin{equation}
  \mathrm{Var}[V] \;=\; (W s)^2 \sum_{k=0}^{\infty} \lambda^{2k}
  \;=\; \frac{(W s)^2}{1-\lambda^2}
  \;=\; \frac{\tau^2}{2\tau-1}\,(W s)^2,
  \label{eq:var}
\end{equation}
%
i.e.\ the membrane variance tracks the \emph{input variance}, scaled by a leak-dependent
constant. Retaining the autocovariance term generalizes \eqref{eq:var} to
%
\begin{equation}
  \mathrm{Var}[V] \;=\; \frac{(W s)^2}{1-\lambda^2}
  \;+\; 2 W^2 \sum_{j=1}^{\infty} \gamma(j)\,\frac{\lambda^{\,j}}{1-\lambda^2},
  \label{eq:var-auto}
\end{equation}
%
which is the key to temporal corruptions: the second term depends on the \emph{time
ordering} of the input, not its marginal distribution. A corruption that preserves $m$ and
$s^2$ but scrambles $\gamma(j)$ moves $\mathrm{Var}[V]$ (and, through higher-order analogues,
the kurtosis and the autocorrelation of $V(t)$) while leaving the naive first-moment picture
untouched. Equations~\eqref{eq:mean}--\eqref{eq:var-auto} are the three levers every
corruption pulls; the rest of the section reads off, corruption by corruption, which lever
moves and therefore which feature detects it.

\subsection{Per-corruption analysis}
\label{sec:theory-percorruption}

\paragraph{Hot-pixel --- DC saturation (predicted AUROC $\approx 1$).}
A stuck/hot pixel injects a persistent constant $\delta$ into $X[t]$ at fixed spatial
locations for the affected channels. Adding a constant $\delta$ to the input in
\eqref{eq:mean} shifts the membrane mean by $\tau W \delta$, an unbounded accumulation that
drives $V$ monotonically toward---and clamps it at---the firing threshold $\theta$. The
affected channels' $\mu_c$ therefore saturate at a value disjoint from the clean range,
producing two non-overlapping clusters in $\varphi$-space. Detection is not merely easy but
\emph{mathematically inevitable}, and it is monotone in severity because larger $\delta$
saturates more channels. This is exactly the measured behaviour (perfect AUROC and Spearman
monotonicity).

\paragraph{Temporal jitter --- autocorrelation destruction (temporal/deep features only).}
Temporal jitter permutes the time order of $X[t]$ within a window. A permutation is
measure-preserving: it leaves the marginal mean $m$ and variance $s^2$---and hence
\eqref{eq:mean} and the \emph{first} term of \eqref{eq:var-auto}---unchanged. What it
destroys is the autocovariance $\gamma(j)$, i.e.\ the second term of \eqref{eq:var-auto} and,
more sharply, the lag-1 autocorrelation of the membrane trace $V(t)$ itself. The leaky filter
\eqref{eq:unroll} is precisely an autocorrelation-generating operation, so jitter is visible
in the \emph{temporal} structure of $V$ even though it is invisible to the static marginal
mean. This predicts, and the data confirm, that jitter is caught by temporal statistics
(lag-1 autocorrelation) and by deeper layers whose larger effective $\tau$ integrate
longer-range temporal structure---the layer-4 slice is the strongest single view for jitter.

\paragraph{Event-rate shift --- global affine drive (single activity scalar).}
An event-rate shift multiplies the global input rate $m \to \alpha m$. By \eqref{eq:mean}
this scales the membrane mean of \emph{every} channel by the same factor $\alpha$, an
approximately affine displacement of $\varphi$ along a single global-activity direction. The
corruption is thus detectable by a one-dimensional projection---a global activity
scalar---which is why a single summary statistic recovers most of the signal, and why the
per-frame AUROC shows a phase transition once $\alpha$ leaves the clean operating band.

\paragraph{Spatial dropout --- variance loss hidden by GAP (anti-detectable under one-sided
scoring).}
Spatial dropout zeros a contiguous spatial \emph{region} of $X[t]$. Within an affected
channel this lowers the spatial input variance and, through \eqref{eq:var}, the membrane
variance $\sigma_c^2$, while leaving the channel mean $\mu_c$ comparatively intact (the
surviving region still drives the neuron). Variance is therefore the single most informative
moment for this corruption---but GAP pools away the \emph{spatial} dimension before the
moments are taken, so the dominant signature (which spatial regions went silent) is
discarded, and what remains makes the membrane look \emph{quieter and more regular}, i.e.\
closer to the clean mean. Under one-sided distance-from-clean scoring this inverts the
detector (AUROC $< 0.5$): the corruption is anti-detectable not because the information is
absent from $V$ but because GAP plus a one-sided rule throws it away. This is the formal
motivation for the spatial-dispersion features (\texttt{phi\_spatial}) and for two-sided
scoring (Section~\ref{sec:method}).

\paragraph{Event flood --- uniform mean shift indistinguishable from a busy scene
(near-chance per frame).}
Event flood adds approximately uniform noise across all pixels, raising $m$ uniformly. By
\eqref{eq:mean} this lifts the mean of \emph{all} channels by nearly the same amount---the
identical signature a genuinely busy, high-activity clean scene produces. Because the clean
distribution already contains such high-activity scenes, the flooded frames are not separable
from the clean tail on the marginal moments alone, giving near-chance per-frame AUROC. The
information that \emph{does} distinguish flood---that the elevated activity is spatially
unstructured (it fills space rather than tracking objects)---again lives in the spatial
dispersion GAP removes, and in the temporal \emph{consistency} of the bias across a sequence.
This predicts the two recoveries we observe: spatial features, and sequence-level
aggregation, which lifts flood to near-perfect AUROC once the per-frame bias is integrated
over a whole recording.

\paragraph{Polarity flip --- invariance by construction ($\approx 0.5$ ceiling, a
membrane-side dead end).}
Polarity flip inverts the ON/OFF sign of the events. The upstream detector was trained to be
approximately polarity-symmetric; to the extent that symmetry holds, $W X$ is near-invariant
to the flip, so $V$---and therefore $\varphi$---barely moves, and the AUROC ceiling is
$\approx 0.5$ \emph{by construction}, not through any deficiency of the detector. The small,
consistent residual below chance (a few points) reflects imperfect symmetry, i.e.\ a faint
directional shift rather than noise. The honest conclusion is that polarity\_flip is largely
undetectable from the membrane and its signal lives on the \emph{input} side (the
positive/negative event-count balance the flip inverts); we scope it accordingly rather than
chase it.

\subsection{Why GAP causes the inversions, and what the theory prescribes}
\label{sec:theory-taxonomy}

The three hard corruptions---spatial\_dropout, event\_flood, polarity\_flip---share a single
mechanism in the analysis above: their discriminative signature lives in structure that
GAP-then-moments discards (spatial dispersion for the first two, input-side sign balance for
the third), while the surviving marginal moments are preserved or even pushed toward the
clean mean. This yields a clean, falsifiable taxonomy that organizes the entire results
section:
%
\begin{itemize}
  \item \textbf{Magnitude corruptions} (hot\_pixel, event\_rate\_shift, event\_flood-as-mean,
        temporal\_jitter via variance) move the membrane moments \emph{away} from clean and
        are caught by distance-based detectors on $\varphi$.
  \item \textbf{Structure corruptions} (spatial\_dropout, and the spatial component of
        event\_flood) leave the marginal moments intact and require either the
        spatial-dispersion features GAP discards or a scoring rule that flags ``too
        typical,'' not only ``too far.''
  \item \textbf{Invariant corruptions} (polarity\_flip) are near-orthogonal to the membrane
        and are an input-side, not a membrane-side, detection problem.
\end{itemize}
%
The theory therefore does three things at once: it \emph{explains} the easy wins as
inevitable, it \emph{predicts} the inversions from first principles, and it \emph{prescribes}
the architecture that addresses them---the spatial branch, two-sided scoring, and sequence
aggregation that the MDD detector (Section~\ref{sec:method}) combines. Crucially, it does so
without any free parameters fit to the corruption labels: every prediction follows from the
PLIF state equation and the definition of the feature.

\section{The Neuromorphic Hardware Advantage}
\label{sec:hardware}

The empirical case for membrane-statistics OOD detection is incomplete without its central
positioning claim, which is not about accuracy at all. Existing OOD detectors are designed
for artificial neural networks executed on von Neumann hardware, where logits, gradients, and
repeated forward passes are cheap. Spiking networks are deployed for the opposite reason---to
run within the severe energy and compute budget of neuromorphic silicon---and in that regime
the standard OOD toolkit is not merely expensive but \emph{architecturally infeasible}.
Membrane statistics are the rare signal that is free precisely where every alternative is
impossible. This reframes the contribution from ``a competitive OOD score'' to ``the only OOD
method that runs on the hardware the model targets.''

\subsection{The membrane potential is a native, already-maintained register}
\label{sec:hardware-register}

On the major neuromorphic platforms---Intel Loihi \cite{davies2018loihi}, the BrainScaleS
analog system \cite{schemmel2010brainscales}, and SpiNNaker \cite{furber2014spinnaker}---each
neuron's membrane potential $V$ is \emph{physical state the substrate already maintains}. It
is not an intermediate that must be recomputed for our purposes; it is the quantity the neuron
updates on every time step in order to decide whether to spike, held in a local register
(digital cores) or a capacitor voltage (analog). Producing our descriptor $\varphi$ from it
requires only a read of that state and a reduction to three per-channel moments---an
$O(\text{channels})$ operation with no additional synaptic integration. In the cost accounting
that matters on-chip---multiply--accumulate operations and memory traffic for weights---%
$\varphi$ adds \emph{zero MACs and zero weight fetches}: it is a byproduct of the forward pass
that has to happen anyway to emit spikes. This is the precise sense in which the method is
``zero additional compute,'' and it is a property of \emph{where} the signal comes from, not
an optimization we apply afterward.

\subsection{Why competing OOD methods cannot run on the target hardware}
\label{sec:hardware-table}

The dominant OOD families each presuppose a capability that the neuromorphic execution model
does not provide. We enumerate them not to claim our accuracy is higher, but to establish that
the comparison cannot even take place on-chip.

\begin{table}[t]
  \centering
  \small
  \begin{tabularx}{\linewidth}{@{}l X X@{}}
    \toprule
    Method & Required extra computation & Why infeasible on neuromorphic HW \\
    \midrule
    \textbf{Vmem-$\varphi$ (ours)} & native $V$ register read + moments & --- (it is free) \\
    MSP \cite{hendrycks2017msp} / Energy \cite{liu2020energy}
      & softmax / log-sum-exp over logits
      & cores emit spikes, not float logits; no on-chip softmax \\
    ODIN \cite{liang2018odin}
      & temperature softmax + input perturbation
      & perturbation needs an input-gradient pass the chip cannot take \\
    ReAct \cite{sun2021react}
      & activation clipping + second forward pass
      & re-forwarding violates the SNN energy budget \\
    ViM \cite{wang2022vim}
      & SVD / dense linear algebra on features
      & no dense float linear-algebra unit; data movement alone too costly \\
    GradNorm \cite{huang2021gradnorm}
      & gradient backpropagation at inference
      & neuromorphic inference has no backward pass \\
    kNN / Mahalanobis on \emph{ANN} features \cite{lee2018mahalanobis,sun2022knn}
      & an auxiliary ANN forward pass
      & requires a second, non-spiking network---defeating the SNN premise \\
    \bottomrule
  \end{tabularx}
  \caption{Standard OOD detectors each require a capability absent from the spiking,
  feed-forward, integer-event execution model; Vmem-$\varphi$ requires none because the signal
  is the neuron state itself.}
  \label{tab:hardware}
\end{table}

The pattern is uniform: every alternative needs either a floating-point readout layer
(softmax, SVD), a second pass (ReAct, ANN-feature methods), or a gradient (ODIN,
GradNorm)---none of which exist in the spiking, feed-forward, integer-event execution model.
Vmem-$\varphi$ needs none of them because the signal is the neuron state itself.

\subsection{The reframed contribution}
\label{sec:hardware-claim}

It follows that the benchmark of Section~\ref{sec:experiments} should be read not as
``$\varphi$ narrowly beats ANN OOD baselines on a GPU'' but as ``$\varphi$ provides usable OOD
detection \emph{inside the compute envelope of the deployment hardware}, where the listed
alternatives cannot be evaluated at all.'' This is a categorical claim about feasibility, and
it is what makes the work a distinct contribution rather than another row in an OOD
leaderboard: for spiking networks on neuromorphic silicon, membrane-statistics OOD detection
is not the best option among many---under a strict on-chip compute budget it is, to our
knowledge, the only option that is free.

\subsection{Scope and honesty caveat}
\label{sec:hardware-caveat}

We measure $\varphi$ in simulation (SpikingJelly on a GPU), not on neuromorphic silicon. Our
cost claim is therefore \emph{architectural}---it follows from $V$ being a register the
substrate already maintains and from $\varphi$ adding no synaptic operations---rather than a
wall-clock or joule measurement on-chip. We state this explicitly: the contribution is that
the signal is free \emph{by construction} on the target hardware, and an on-silicon energy
measurement is identified as future work. We likewise restrict the infeasibility argument to
the standard execution models of current neuromorphic platforms; co-processor or hybrid
configurations could in principle host some of the listed methods at a cost, which only
sharpens the comparison rather than changing its direction.

% --- Standalone bibliography (skipped when embedded). Manual thebibliography so
%     the standalone build needs no .bib file and no bibtex pass. Replace with a
%     real \bibliography{...} when merging into the full manuscript.
\ifdefined\theoryhardwareEMBED\else
\begin{thebibliography}{99}

\bibitem{fang2021plif}
W.~Fang, Z.~Yu, Y.~Chen, T.~Masquelier, T.~Huang, and Y.~Tian.
\newblock Incorporating learnable membrane time constant to enhance learning of
  spiking neural networks.
\newblock In \emph{ICCV}, 2021.

\bibitem{davies2018loihi}
M.~Davies et~al.
\newblock Loihi: A neuromorphic manycore processor with on-chip learning.
\newblock \emph{IEEE Micro}, 38(1):82--99, 2018.

\bibitem{schemmel2010brainscales}
J.~Schemmel, D.~Br\"uderle, A.~Gr\"ubl, M.~Hock, K.~Meier, and S.~Millner.
\newblock A wafer-scale neuromorphic hardware system for large-scale neural
  modeling.
\newblock In \emph{ISCAS}, pages 1947--1950, 2010.

\bibitem{furber2014spinnaker}
S.~B. Furber, F.~Galluppi, S.~Temple, and L.~A. Plana.
\newblock The SpiNNaker project.
\newblock \emph{Proceedings of the IEEE}, 102(5):652--665, 2014.

\bibitem{hendrycks2017msp}
D.~Hendrycks and K.~Gimpel.
\newblock A baseline for detecting misclassified and out-of-distribution
  examples in neural networks.
\newblock In \emph{ICLR}, 2017.

\bibitem{liu2020energy}
W.~Liu, X.~Wang, J.~D. Owens, and Y.~Li.
\newblock Energy-based out-of-distribution detection.
\newblock In \emph{NeurIPS}, 2020.

\bibitem{liang2018odin}
S.~Liang, Y.~Li, and R.~Srikant.
\newblock Enhancing the reliability of out-of-distribution image detection in
  neural networks.
\newblock In \emph{ICLR}, 2018.

\bibitem{sun2021react}
Y.~Sun, C.~Guo, and Y.~Li.
\newblock ReAct: Out-of-distribution detection with rectified activations.
\newblock In \emph{NeurIPS}, 2021.

\bibitem{wang2022vim}
H.~Wang, Z.~Li, L.~Feng, and W.~Zhang.
\newblock ViM: Out-of-distribution with virtual-logit matching.
\newblock In \emph{CVPR}, 2022.

\bibitem{huang2021gradnorm}
R.~Huang, A.~Geng, and Y.~Li.
\newblock On the importance of gradients for detecting distributional shifts in
  the wild.
\newblock In \emph{NeurIPS}, 2021.

\bibitem{lee2018mahalanobis}
K.~Lee, K.~Lee, H.~Lee, and J.~Shin.
\newblock A simple unified framework for detecting out-of-distribution samples
  and adversarial attacks.
\newblock In \emph{NeurIPS}, 2018.

\bibitem{sun2022knn}
Y.~Sun, Y.~Ming, X.~Zhu, and Y.~Li.
\newblock Out-of-distribution detection with deep nearest neighbors.
\newblock In \emph{ICML}, 2022.

\end{thebibliography}
\end{document}
\fi
